%\documentclass{amsart}
\documentclass[leqno]{amsart}
\usepackage{amssymb}
\usepackage{amsmath,amsfonts, enumerate}
\usepackage{color}
\usepackage{amsthm}
\usepackage[backref]{hyperref}%hyperref生成超链接,backref可看文献引用多少次
\usepackage{mathrsfs}
\usepackage{txfonts}
\usepackage{cleveref}%(引用自动生成定理,引理字样的宏包)


 \usepackage{appendix}
\usepackage{showkeys}

\begin{document}

\newtheorem{theorem}{Theorem}    %[section]
\newtheorem{proposition}[theorem]{Proposition}
\newtheorem{conjecture}[theorem]{Conjecture}
\newtheorem{corollary}[theorem]{Corollary}
\newtheorem{lemma}[theorem]{Lemma}
\newtheorem{sublemma}[theorem]{Sublemma}
\newtheorem{observation}[theorem]{Observation}
\newtheorem{remark}[theorem]{Remark}
\newtheorem{definition}[theorem]{Definition}
\newtheorem{notation}[theorem]{Notation}
\newtheorem{question}[theorem]{Question}
\newtheorem{questions}[theorem]{Questions}
\newtheorem{example}[theorem]{Example}
\newtheorem{problem}[theorem]{Problem}
\newtheorem{exercise}[theorem]{Exercise}

\numberwithin{theorem}{section}
\numberwithin{theorem}{section}
\numberwithin{equation}{section}

\newcommand{\ri}{{\rm{i}}}
\newcommand{\inner}[1]{ \langle  #1 \rangle}
\newcommand{\Norm}[1]{ \left\|  #1 \right\| }
\newcommand{\set}[1]{ \left\{ #1 \right\} }
\newcommand{\modulo}[2]{[#1]_{#2}}
\newcommand{\bra}[1]{\Bigg(  #1 \Bigg)}
\newcommand{\cir}[1]{\left(  #1 \right)}
\newcommand{\wh}[1]{\widehat{#1}}
\newcommand{\wt}[1]{\widetilde{#1}}
\newcommand{\cl}[1]{{\rm  #1 }}
\newcommand{\lp}[1]{L^{#1}(\RR^{n})}
\newcommand{\np}[1]{L_{#1}(\mathcal{M})}
\newcommand{\nlp}[1]{L_{#1}(\RR^{d};\A_{\theta}^{d})}


\def\one{\mathbf 1}
\def\sn{{{S}^{n-1}}}
\def\qd{\,{\mathchar'26\mkern-12mu d}}
\def\p{\partial}
\def\rp{ ^{-1} }
\def\supp{{\rm supp}}
\def\ted{\mathbb{T}_{\theta}^{d}}
\def\td{\mathbb{T}^{d}}
\def\zd{\mathbb{Z}^{d}}
\def\rd{\mathbb{R}^{d}}
\def\cci{C_{c}^{\infty}(\rd)}
\def\e{\mathbf e}
\def\BMO{{\rm BMO}}
\def\bmo{{\rm bmo}}
\def\Tr{{\rm Tr}}


\def\MM{\mathbb{M}}
\def\DD{\mathbb{D}}
\def\NN{\mathbb{N}}
\def\QQ{{\mathbb{Q}}}
\def\SS{\mathbb{S}}
\def\RR{{\mathbb R}}
\def\TT{{\mathbb T}}
\def\HH{{\mathbb H}}
\def\ZZ{{\mathbb Z}}
\def\CC{{\mathbb C}}
\def\EE{\mathbb E}
\def\GG{\mathbb G}

\def\htv{\mathbf V}
\def\htd{\mathbf D}
\def\htb{\mathbf B}


\def\H{{\mathcal H}}
\def\L{\mathcal{L}}
\def\S{\mathcal{S}}
\def\F{{\mathcal F}}
\def\G{{\mathcal G}}
\def\M{{\mathcal M}}
\def\B{{\mathcal B}}
\def\C{{\mathcal C}}
\def\T{{\mathcal T}}
\def\I{{\mathcal I}}
\def\E{{\mathcal E}}
\def\V{{\mathcal V}}
\def\W{{\mathcal W}}
\def\U{{\mathcal U}}
\def\A{{\mathcal A}}
\def\R{{\mathcal R}}
\def\O{{\mathcal O}}
\def\D{{\mathcal D}}
\def\N{{\mathcal N}}
\def\P{{\mathcal P}}
\def\K{{\mathcal K}}


\title[$L\log L$-estimate of Commutator]{P\'{e}rez's $L\log L$-estimate of commutator beyond doubling measure}

\thanks{{\it 2000 Mathematics Subject Classification:} Primary:  Secondary: }

\thanks{{\it Key words:} Non-doubling measure, dyadic Hilbert transform, commutator, singular integral operator}


\author{Fedor Sukochev}
\address{School of Mathematics and Statistics, UNSW, Kensington, NSW 2052, Australia}
\email{f.sukochev@unsw.edu.au}



\author{Fulin Yang}
\address{School of Mathematics, Harbin Institute of Technology, Harbin 150001, China}
\email{fulinyoung@hit.edu.cn}


\author{Dmitriy Zanin}
\address{School of Mathematics and Statistics, UNSW, Kensington, NSW 2052, Australia}
\email{d.zanin@unsw.edu.au}


\author{Dejian Zhou}
\address{}
\email{}

%\markboth{F. A. Sukochev, F. L. Yang and D. Zanin}
%{Gagliardo-Nirenberg inequality}





\date{}


\begin{abstract}
	In this paper, we investigate  P\'{e}rez's $L\log L$-estimate in the setting of non-doubling measure. Let $\mu$ be a $m$-increasing positive Borel measure on $\RR$ and let $H_{\D}$ be the associated dyadic Hilbert transform from dyadic cubes system $\D$ on $\RR$. For each $\lambda>0$, we show that the following $L\log L$-estimate holds:
	\begin{align*}
		\mu(\set{|[M_{\varphi},H_{\D}]f|>\lambda})\leq C
		\int_{\RR}\frac{\Norm{\varphi}_{\BMO_{\mu,2}}|f(y)|}{\lambda}\left(1+\log^{+}(\frac{\Norm{\varphi}_{\BMO_{\mu,2}}|f(y)|}{\lambda})\right)d\mu(y)
	\end{align*}
	As a byproduct, we also show the $L_p$-boundedness of $[M_{\varphi},H_{\D}]$ for $1<p\leq2$; however, unlike the classical case, this commutator is not bounded on $L_p(\mu)$ whenever $2<p<\infty$, $\mu$ is non-doubling.
\end{abstract}

\maketitle













%%%%%%%%%%%%%%%%%%%%%%%%%%%%%%%%%%%%%%%%%%%%%%%%%%%%%%%%%%%%%%%%%%%%%%%
\section{Introduction}
On Euclidean space, singular integral and its commutator play an important role in modern harmonic analysis, functional analysis and PDEs et. The Calder\'{o}n--Zygmund singular integral operator is defined by setting
$$Tf(x)={\rm p.v.} \int_{\RR^{d}}K(x-y)f(y)dy,\quad x\in \mathbb{R}^d, f\in\S(\RR^{d}),$$
where the kernel function $K:\RR^{d}\setminus\{0\}\rightarrow\CC$ is measurable and satisfies the assumptions
\begin{enumerate}[\rm(1)]
\item $|K(x)|\leq C\frac{1}{|x|^{d}}$ for $x\neq0$ and some $C>0$; 
\item\label{kernel smoothness} $\left|K(x-y)-K(x)\right|\leq C\frac{|y|}{|x|^{d+1}}$ for $|x|\geq2|y|>0$ and some constant $C>0$.
\end{enumerate}
It is well-known that $T$ is of strong type $(p,p)$, $1<p<\infty$, and  of weak type $(1,1)$ under the assumption that $T$ is bounded on $L_{2}(\RR^{d})$ (see \cite{Gr-book1}).
For a locally integrable function $\varphi$ on $\RR^{d}$, the commutator $[M_{\varphi},T]$ is given by $$[M_{\varphi},T]=M_{\varphi}T-TM_{\varphi}$$
where $M_{\varphi}$ is the multiplication operator on $L_{2}(\RR^{d})$. In the remarkable paper \cite{CRW1976}, Coifman-Rochberg-Weiss proved that the operator $[M_{\varphi},T]$ is bounded on $L_{p}(\RR^{d})$, $1<p<\infty$, if and only if $\varphi$ is a $\BMO(\RR^{d})$ function. However, unlike the singular integral operator, the commutator $[M_{\varphi},T]$ fails to be of weak $(1,1)$ provided $\varphi\in\BMO(\RR^{d})$. This is interesting and unexpected. One can find an example in P\'{e}rez's paper \cite[Section 5]{Pe1995}. Moreover, in the same paper, P\'{e}rez established the $L\log^{+}L$ estimate of $[M_{\varphi},T]$ but it should be in the following form. 

\begin{theorem}[{\cite[Theorem 1.1]{Pe1995} }]\label{Perez thm}
	Let $T$ be a Calder\'{o}n--Zygmund singular integral operator and let $\varphi\in\BMO(\RR^{d})$. Then, for all $\lambda>0$, there exists a positive constant $C$, independent of $\lambda$, $\varphi$ and $f$, such that
	\begin{align}\label{LlogL correct form}
		\left|\set{y\in\RR^{d}:\left|[M_{\varphi},T]f(y)\right|>\lambda}\right|\leq C\int_{\RR^{d}}\frac{\Norm{\varphi}_{\BMO}|f(y)|}{\lambda}\left(1+\log^{+}(\frac{\Norm{\varphi}_{\BMO}|f(y)|}{\lambda})\right)dy.
	\end{align}
\end{theorem}

 We definitely convince that the following fact is known to experts; see e.g.\cite{CF2003,CF2007, PP2001}. We still include it for reader's convenience.
% We hope it to be clear to everyone so that we write down our comment here.
%The following comment maybe known to experts, we still include it for reader's convenience.
\begin{remark}
Recall that, in the initial version of Theorem \ref{Perez thm} in \cite{Pe1995}, \eqref{LlogL correct form} was written as
	\begin{align}\label{false result}
		\left|\set{y\in\RR^{d}:\left|[M_{\varphi},T]f(y)\right|>\lambda}\right|\leq C\Norm{\varphi}_{\BMO}\int_{\RR^{d}}\frac{|f(y)|}{\lambda}\left(1+\log^{+}(\frac{|f(y)|}{\lambda})\right)dy.
	\end{align}
However, this assertion is incorrect. Indeed, let us substitute $\varphi$ and $\lambda$ as $t\varphi$ and $t\lambda$ on both sides of \eqref{false result} for $t>0$. Note that, for $t>0$,
\begin{align*}
	\left|\set{y\in\RR^{d}:\left|[M_{t\varphi},T]f(y)\right|>t\lambda}\right|
	=\left|\set{y\in\RR^{d}:\left|[M_{\varphi},T]f(y)\right|>\lambda}\right|.
\end{align*}
It follows that
\begin{align*}
	\left|\set{y\in\RR^{d}:\left|[M_{\varphi},T]f(y)\right|>\lambda}\right|\leq C\Norm{t\varphi}_{\BMO}\int_{\RR^{d}}\frac{|f(y)|}{t\lambda}\left(1+\log^{+}(\frac{|f(y)|}{t\lambda})\right)dy\\
	=C\Norm{\varphi}_{\BMO}\int_{\RR^{d}}\frac{|f(y)|}{\lambda}\left(1+\log^{+}(\frac{|f(y)|}{t\lambda})\right)dy.
\end{align*}
Since the factor on the left hand side is independent of $t$, letting $t\to+\infty$ on the right hand side, it follows that
\begin{align*}
	\left|\set{y\in\RR^{d}:\left|[M_{\varphi},T]f(y)\right|>\lambda}\right|\leq C\Norm{\varphi}_{\BMO}\frac{\Norm{f}_{L_{1}(\RR^{d})}}{\lambda}.
\end{align*}
This implies that $[M_{\varphi},T]$ is of weak type $(1,1)$ and it contradicts to the known fact given in \cite[Section 5]{Pe1995}. 
\end{remark}

In the last three decades, there are lots of literature concern related topic, here we only mention some recent advance of them: \cite{FSZ2024,HLW2017,Hy2021,LXY2024,Wi2020}. In the present paper, we mainly concern the $L\log L$-estimate of the commutator of the dyadic Hilbert transform on $\RR$, however, with respect to some non-doubling measure $\mu$.



Let $\D$ be the standard dyadic interval system on $\RR$, and let $\mu$ be a Borel measure on $\RR$ with $0<\mu(Q)<\infty$ as $Q\in\D$. For our purpose, we always assume that $\mu$ is a complete $\sigma$-additive measure with $\mu(\RR)=\infty$ in the sequel. The (weak) Lebesgue spaces with respect to the measure $\mu$ are denoted by $L_p(\mu)$, $L_{p,\infty}(\mu)$, $0<p<\infty$. Given $Q\in\D$, write $Q_{-}$ and $Q_{+}$ as the (left and right) dyadic children of $Q$, and denote $\wh{Q}$ the dyadic parent of $Q$.  
Set
\begin{align*}
h_{Q}=m(Q)^{\frac{1}{2}}\left(\frac{\chi_{Q_{-}}}{\mu(Q_{-})}-\frac{\chi_{Q_{+}}}{\mu(Q_{+})}\right)
\quad{\rm with}\quad m(Q):=\frac{\mu(Q_{-})\mu(Q_{+})}{\mu(Q)}.
\end{align*}
The system $\{h_{Q}\}_{Q\in\D}$ is an orthonormal basis in $L_{2}(\mu)$. 
In addition, for each $Q\in\D$, we have 
\begin{align}\label{hQ norm}
\Norm{h_{Q}}_{L_{1}(\mu)}=2m(Q)^{\frac{1}{2}},\quad\Norm{h_{Q}}_{L_{\infty}(\mu)}\simeq m(Q)^{-\frac{1}{2}}.
\end{align}
Following \cite[page 50]{LMP2014}, we define the dyadic Hilbert transform as below
\begin{align*}
H_{\D}f(x)=\sum_{Q\in\D}\inner{f,h_{Q}}(h_{Q_{-}}(x)-h_{Q_{+}}(x))
=\sum_{Q\in\D}\sigma(Q)\inner{f,h_{\wh{Q}}}h_{Q}(x),
\end{align*}
where $f\in L_{2}(\mu)$ and $\sigma(Q)=1$ if $Q=\wh{Q}_{-}$ and $\sigma(Q)=1$ if $Q=\wh{Q}_{+}$. 
This dyadic Hilbert transform is a generalization of the classical Haar shift operator on $\RR$ with respect to Lebesgue measure which is bounded on $L_{p}(\RR)$ $(1<p<\infty)$ and bounded from $L_{1}(\RR)$ to $L_{1,\infty}(\RR)$ (see \cite{CMP2012}).
The $L_{2}(\mu)$-adjoint operator of $H_{\D}$ can be written as $$H_{\D}^{*}=\sum_{Q\in\D}\sigma(Q)\inner{f,h_{Q}}h_{\wh{Q}}(x).$$

The following definition was firstly introduced in \cite[page 50]{LMP2014} (in what follows, unexplained terminology is referred to Section \ref{sec-pre}). 

\begin{definition}
We say that $\mu$ is $m$-increasing if there exists constant $0<C<\infty$ such that
$$m(Q)\leq Cm(\wh{Q}),\quad \forall  Q\in\D.$$
We say that $\mu$ is $m$-decreasing if there exists constant $0<C<\infty$ such that
$$m(\wh{Q})\leq C m(Q),\quad \forall Q\in\D.$$
\end{definition}
In fact, these two definitions characterize the equivalent condition for the boundedness of $H_{\D}$ and $H_{\D}^{*}$ (see \cite[Theorem 2.5]{LMP2014}).
More precisely, they proved the following equivalence:
\begin{enumerate}[{\rm (i)}]
	\item The Borel measure $\mu$ is $m$-increasing;
	\item $H_{\D}$ is bounded from $L_{1}(\mu)$ to $L_{1,\infty}(\mu)$;
	\item $H_{\D}$ is bounded on $L_{p}(\mu)$ for $1<p<2$;
	\item  $H_{\D}^{*}$ is bounded on $L_{q}(\mu)$ for $2<q<\infty$.
\end{enumerate}
Noting that $H_{\D}$ and $H_{\D}^{*}$ are bounded on $L_{2}(\mu)$. The $L_{2}(\mu)$-boundedness is independent of the property $m$-increasing or $m$-decreasing. With $L_{p}(\mu)$-boundedness of $H_{\D}$ in hand, it is natural to investigate the boundedness of the commutator of $H_{\D}$.
Recently, Bonami et al. \cite{BJXYZ2023} studied $H_1-L_1$ estimate of the commutator $[M_{\varphi}, H_{\mathcal{D}}]$ via bilinear decomposition of martingale Hardy and BMO spaces; Borges et al. \cite{BCPW2024} investigated the $L_p$-estimate with $1<p<\infty$ of such commutator of another type dyadic Hilbert transform in the setting of "sibling balanced" measure. Hence, it is natural to continue to consider the $L\log L$-estimate. 



\begin{theorem}\label{endpoint estimate}
Let $\mu$ be a $m$-increasing Borel measure on $\RR$ with $0<\mu(Q)<\infty$ for each $Q\in\D$.  For all $\lambda>0$, if $\varphi\in\BMO_{\mu,2}$ and $f\in L\log^{+}L(\mu)$, then there exists a positive constant $C$, independent of $\lambda$, $\varphi$ and $f$, such that
\begin{align*}
\mu(\set{|[M_{\varphi},H_{\D}]f|>\lambda})\leq C
\int_{\RR}\frac{\Norm{\varphi}_{\BMO_{\mu,2}}|f(y)|}{\lambda}\left(1+\log^{+}(\frac{\Norm{\varphi}_{\BMO_{\mu,2}}|f(y)|}{\lambda})\right)d\mu(y).
\end{align*}
\end{theorem}
The main difficulty to prove this theorem is how to treat the non-doubling condition of $\mu$. For a cube $Q\in\D$, comparing $\mu(\wh{Q})$ with $\mu(Q)$ should be deal carefully. In order to avoid comparing them, we use the new Calder\'{o}n-Zygmund decomposition (established in \cite{LMP2014}) with respect to the non-doubling measure instead of the classical one. 
%The classical Calder\'{o}n-Zygmund decomposition decompose the desired function to good and bad functions but 
This new decomposition decompose the desired function to three parts involving the parents of the selected cubes so that comparing $\mu(\wh{Q})$ with $\mu(Q)$ transfers to comparing $m(\wh{Q})$ with $\mu(Q)$. When $\mu$ is a $m$-increasing Borel measure, we know $m(\wh{Q})\leq\mu(Q)$ and it works. 

 

It is noteworthy that the martingale $\BMO_{\mu,2}$ space is different with the classical one, saying $\bmo_{\mu,2}$ (see Section 2.2), in Euclidean space with respect to general $\mu$. The $\BMO_{\mu,2}$ is defined by the oscillation with $f_{\widehat{Q}}$ instead of $f_{Q}$. When $\mu$ is doubling, $\BMO_{\mu,2}$ and $\bmo_{\mu,2}$ are isomorphic with equivalent norm.  For $d=1$ and Lebesgue measure $\mu$ on the line, \eqref{LlogL correct form} holds for $\varphi$ belonging to $\bmo_{\mu,2}$ and some $\BMO$ space by translation of $\bmo_{\mu,2}$ (see \cite{Me2003}).
In the setting of non-doubling $\mu$, $\BMO_{\mu,2}$ is a subspace of $\bmo_{\mu,2}$ and the commutator $[M_{\varphi},H_{\D}]$ cannot be controlled only by $\Norm{\varphi}_{\bmo_{\mu,2}}$. This requires us treat $\BMO_{\mu,2}$ and $\bmo_{\mu,2}$ delicately. In addition, another key point to use $\BMO_{\mu,2}$ instead of $\bmo_{\mu,2}$ is the John-Nirenberg inequality valid in $\BMO_{\mu,2}$ (see Subsection 2.2).

Our second result is the $L_{p}$-boundedness of $[M_{\varphi},H_{\D}]$ with respect to a $m$-increasing Borel measure on $\RR$. It is proved with the help of Theorem \ref{endpoint estimate}.

\begin{theorem}\label{Lp estimate}
Let $\mu$ be a $m$-increasing Borel measure on $\RR$ with $0<\mu(Q)<\infty$ for each $Q\in\D$. If $\varphi\in\BMO_{\mu,2}$, then $[M_{\varphi},H_{\D}]$ is bounded on $L_{p}(\mu)$  for  $1<p\leq2$ and $[M_{\varphi},H_{\D}^{*}]$ is bounded on $L_{q}(\mu)$  for  $2\leq q<\infty$. 
%\begin{enumerate}[\rm(a)]
%\item $[M_{\varphi},H_{\D}]$ is bounded on $L_{p}(\mu)$  for  $1<p\leq2$;
%\item $[M_{\varphi},H_{\D}^{*}]$ is bounded on $L_{q}(\mu)$  for  $2\leq q<\infty$. 
%\end{enumerate}
More precisely,
\begin{align*}
\Norm{[M_{\varphi},H_{\D}]}_{L_{p}(\mu)\rightarrow L_{p}(\mu)}
\leq C\Norm{\varphi}_{\BMO_{\mu,2}}\;\mbox{and}\;
\Norm{[M_{\varphi},H_{\D}^{*}]}_{L_{q}(\mu)\rightarrow L_{q}(\mu)}
\leq C\Norm{\varphi}_{\BMO_{\mu,2}},
\end{align*}
for some positive constant $C$ independent of $\varphi$.
\end{theorem}

In the setting of non-doubling measure, we only have the $L_{p}$-boundedness of $[M_{\varphi},H_{\D}]$ for $1<p\leq2$ instead of $1<p<\infty$. When $q>2$, $[M_{\varphi},H_{\D}]$ may be unbounded since there are some non-doubling measures such that $H_{\D}$ is unbounded on $L_{q}(\mu)$ (see \cite[Section 4.1]{LMP2014}). In fact, in the proof of this theorem, we need to deal with the paraproduct and its adjoint. The paraproduct is bounded on $L_{p}(\mu)$ with $1<p\leq2$ for a general positive Borel measure (see \Cref{remark on paraproduct}) and therefore its adjoint is bounded on $L_{q}(\mu)$ for $2\leq q<\infty$. Maybe this is the key point to limit the $L_{p}$-boundedness of $[M_{\varphi},H_{\D}]$ to $1<p\leq2$.
Besides, if one further assume that the measure $\mu$ is ``sibling balanced", then the $L_{p}$-boundedness of $[M_{\varphi},H_{\D}]$ holds for all $1<p<\infty$; see \cite{BCPW2024} for more details. In the same paper \cite[Theorem A]{BCPW2024}, the authors also showed that the following estimate generally fails
$$\|\varphi\|_{\BMO_{\mu,2}}\lesssim \|[M_{\varphi},\mathcal{H}_{\mathcal{D}}]\|_{L_2(\mu)\to L_2(\mu)}.$$
Thus, we do not consider the converse direction of \Cref{endpoint estimate} and \Cref{Lp estimate}.

As the argument in previous, when $\mu$ is the Lebesgue measure on $\RR$, $\BMO_{\mu,2}$ and $\bmo_{\mu,2}$ are isomorphic with equivalent norm. In this case, \Cref{endpoint estimate} and \Cref{Lp estimate} give the following corollary.

\begin{corollary}
When $\mu$ is the Lebesgue measure, $H_{\D}$ is the classical dyadic Hilbert transform on $\RR$. 
Denote $\bmo_{2}(\RR)$ the $\bmo_{\mu,2}$ with $\mu$ as the Lebesgue measure on $\RR$. 
In this case, for $\varphi\in\bmo_{2}(\RR)$, $[M_{\varphi},H_{\D}]$ is bounded on $L_{p}(\RR)$  for  $1<p<\infty$ and $[M_{\varphi},H_{\D}]$ admits the $L\log^{+}L(\RR)$ estimate as in \Cref{endpoint estimate}.
\end{corollary}
%\begin{center}
%$\mu$ is $m$-decreasing 
%$\Longleftrightarrow$ $H_{\D}^{*}$ is bounded from $L_{1}(\mu)$ to $L_{1,\infty}(\mu)$ 
%$\Longleftrightarrow$ $H_{\D}^{*}$ is bounded on $L_{p}(\mu)$ for $1<p<2$ 
%$\Longleftrightarrow$ $H_{\D}$ is bounded on $L_{q}(\mu)$ for $2<q<\infty$.
%\end{center}


This paper is organized as follows. In Section \ref{sec-pre}, we introduce a  Calder\'{o}n-Zygmund decomposition associated to non-doubling measure, the definitions of martingale $\BMO$ spaces and several properties of some symmetric function spaces. In Section \ref{sec-Lpe}, we provide the proof of \Cref{Lp estimate}. Note that we actually prove \Cref{Lp estimate} via Theorem \ref{endpoint estimate}. Finally, in Section \ref{sec-end}, we show \Cref{endpoint estimate}. 

Throughout the paper, we use $A\lesssim B$ to denote the statement that $A\leq CB$ for some constant $C>0$ independent of $A,B$.


%Let $L\log^{+}L(\mu)$ be the $L\log^{+}L$ function space on $\RR$ with respect to measure $\mu$ and denote $\BMO_{\mu,2}$ the function space on $\RR$ with respect to measure $\mu$ involving the locally integrable functions such as
%\begin{align*}
%	\Norm{f}_{\BMO_{\mu,2}}:=
%	\sup_{n}\sup_{Q\in\D_{n}}\left(\frac{1}{\mu(Q)}\int_{Q}|f-f_{\widehat{Q}}|^{2}d\mu\right)^{\frac{1}{2}}
%	<\infty
%\end{align*}
%where $\wh{Q}$ is the father of $Q\in\D$ and $f_{Q}$ denotes the $\mu$-average of $f$ over cube $Q\in\D$.
%Our first result is the following $L\log L$ estimate of $[M_{\varphi},H_{\D}]$ with respect to a $m$-increasing Borel measure on $\RR$.
%
%where $L_{2}(\mu)$ is the square integrable function space on $\RR$ with respect to measure $\mu$ (Also, in what follows, denote $L_{1,\infty}(\mu)$ and $L_{p}(\mu)$, $0<p\leq\infty$, the weak $(1,1)$ and the $L_{p}$ function space on $\RR$ with respect to measure $\mu$). 

\section{Preliminaries}\label{sec-pre}
Basic material used in the paper is collected in this section.

\subsection{Calder\'{o}n-Zygmund decomposition}
Throughout the paper, let $\mathcal{D}$ be the standard dyadic interval system on $\RR$. For each $n\in\mathbb{Z}$, write
$$\mathcal{D}_n=\{Q\in \mathcal{D}: \ell(Q)=2^{-n}\},\quad \mathcal{F}_n=\sigma(\mathcal{D}_n),$$
where $\ell(Q)$ denotes the length of $Q$. Then $(\mathcal{F}_n)_{n\in\mathbb{Z}}$ is an increasing filtration. 
Let $\mu$ be a Borel measure on $\RR$ such that $\mu(Q)<\infty$ for $Q\in\D$.
For each $n$, the conditional expectation with respect to $\mathcal{F}_n$ is given by (see e.g. \cite[p.\,2]{Lo1993})
\begin{align}\label{expectation definition}
	\EE_{n}(f)=\sum_{Q\in\D_{n}}\frac{1}{\mu(Q)}\int_{Q}f(y)d\mu(y)\chi_{Q}.
\end{align}
Given $f\in L_1(\mu)$, $(\EE_n(f))_{n\in\mathbb{Z}}$ forms a martingale.  For $f\in L_{1}(\mu)$, we set $f_{Q}=0$ if $\mu(Q)=0$ and define the dyadic maximal function associated to $\mu$ by
$$M_{\D}f(x)=\sup_{\substack{Q\in\D\\Q\ni x}}\frac{1}{\mu(Q)}\int_{Q}|f(x)|d\mu(x).$$
This is actually the well-known Doob maximal function $\sup_n|\EE_n(f)|$; see \cite{Pi2016} or \cite[p.\,33]{Lo1993}. It is of weak type $(1,1)$ and strong type $(p,p)$ with $1\leq p\leq \infty$. We refer the reader to \cite{Ga1973,HNVW2016,Lo1993,Pi2016,We1994} for more background about martingale theory.

When $\mu$ is a non-doubling measure, the classical Calder\'{o}n-Zygmund decomposition is of no use.
This requires us to use the following Calder\'{o}n-Zygmund decomposition in the setting of non-doubling measure.

\begin{lemma}[{\cite[Theorem 2.1]{LMP2014}}]\label{CZ decomposition}
Let $\mu$ be a Borel measure on $\RR$ such that $\mu(Q)<\infty$ for each $Q\in\D$ and $\mu(\RR)=\infty$.
Given $f\in L^{1}(\mu)$ and $\alpha>0$, let $\Omega_{\alpha}=\set{M_{\D}f>\alpha}$. There exists a collection of disjoint maximal dyadic cubes $\{Q_{j}\}_{j}$ covering $\Omega_{\alpha}$ and there is a related decomposition $f=g+b+\beta$ with 
\begin{align*}
&g(x)=f(x)\chi_{\RR^{d}\setminus\Omega_{\alpha}}(x)+\sum_{j}f_{\wh{Q}_{j}}\chi_{Q_{j}}(x)
+\sum_{j}(f_{Q_{j}}-f_{\wh{Q}_{j}})\frac{\mu(Q_{j})}{\mu(\wh{Q}_{j})}\chi_{\wh{Q}_{j}}(x),\\
&b(x)=\sum_{j}b_{j}(x)=(f(x)-f_{Q_{j}})\chi_{Q_{j}}(x),\\
&\beta(x)=\sum_{j}\beta_{j}(x)=\sum_{j}(f_{Q_{j}}-f_{\wh{Q}_{j}})
\left(\chi_{Q_{j}}(x)-\frac{\mu(Q_{j})}{\mu(\wh{Q}_{j})}\chi_{\wh{Q}_{j}}(x)\right),
\end{align*}
where the maximality of $\{Q_{j}\}_{j}$ is understood in the sense that, for all dyadic intervals $Q\varsupsetneq Q_{j}$, we have $|f|_{Q}\leq\alpha\leq|f|_{Q_{j}}$. 
Additionally, they satisfy the following properties:
\begin{enumerate}[\rm(i)]
\item\label{Omega upper bd} $\mu(\Omega_{\alpha})\lesssim\frac{\Norm{f}_{L_{1}(\mu)}}{\alpha}$;
%\item\label{f out Omega bd} $|f(x)|\leq\alpha$, a.e. $x\in\RR^{n}\setminus\Omega_{\alpha}$;
\item\label{g upper bd} $\Norm{g}_{L_{p}(\mu)}\leq C_{p}\alpha^{p-1}\Norm{f}_{L_{1}(\mu)}$ for all $1\leq p<\infty$;
\item\label{b cancellation} $\supp(b_{j})\subset Q_{j}$, $\displaystyle\int_{\RR^d}b_{j}(x)d\mu(x)=0$ and $\displaystyle\sum_{j}\Norm{b_{j}}_{L_{1}(\mu)}\leq2\Norm{f}_{L_{1}(\mu)}$;
\item\label{beta cancellation} $\supp(\beta_{j})\subset \wh{Q}_{j}$, $\displaystyle\int_{\RR^d}\beta_{j}(x)d\mu(x)=0$ and $\displaystyle\sum_{j}\Norm{\beta_{j}}_{L_{1}(\mu)}\leq4\Norm{f}_{L_{1}(\mu)}$.
\end{enumerate}
\end{lemma}

This lemma is closely related to Gundy's martingale decomposition \cite{Gu1968} and also can be discerned in \cite{CMP-book}. The feature of the maximal cubes in the lemma  plays an important role in characterizing the weak $(1,1)$ of $H_{\D}$ and the $L\log L$ estimate of the commutator $[M_{\varphi},H_{\D}]$ in the setting of non-doubling Borel measure.





\subsection{$\BMO$ space}
For our purpose, let us introduce the definitions and some properties of different martingale $\BMO$ spaces.

\begin{definition}
Let $p\in[1,\infty)$. The  $\BMO_{\mu,p}$ space associated to $\D$ is defined to be the set of all locally integrable functions $f$ defined on $\RR$ with respect to $\mu$ satisfying
\begin{align*}
\Norm{f}_{\BMO_{\mu,p}}:=
\sup_{n}\sup_{Q\in\D_{n}}\left(\frac{1}{\mu(Q)}\int_{Q}|f-f_{\widehat{Q}}|^{p}d\mu\right)^{\frac{1}{p}}
<\infty,
\end{align*}
where $\wh{Q}$ is the father of $Q\in\D$.

The $\bmo_{\mu,p}$ space associated to $\D$ is defined to be the set of all locally integrable functions $f$ defined on $\RR$ with respect to $\mu$ satisfying
\begin{align*}
\Norm{f}_{\bmo_{\mu,p}}:=
\sup_{n}\sup_{Q\in\D_{n}}\left(\frac{1}{\mu(Q)}\int_{Q}|f-f_{Q}|^{p}d\mu\right)^{\frac{1}{p}}<\infty.
\end{align*}

The $\bmo_{\mu,\mathsf{d}}$ space associated to $\D$ is defined to be the set of all locally integrable functions $f$ defined on $\RR$ with respect to $\mu$ satisfying
\begin{align*}
\Norm{f}_{\bmo_{\mu,\mathsf{d}}}:=\sup_{n}\Norm{\EE_{n}f-\EE_{n-1}f}_{L_{\infty}(\mu)}<\infty.
\end{align*}
\end{definition}
It is obvious $$\Norm{f}_{\bmo_{\mu,p}}\leq\Norm{f}_{\BMO_{\mu,p}}\quad\mbox{for}\quad p\in[1,\infty).$$
In the case $p=2$, it is clear that 
\begin{align}\label{BMO intersection space}
\BMO_{\mu,2}=\bmo_{\mu,2}\bigcap\bmo_{\mu,\mathsf{d}},
\end{align}
with norm $$\Norm{f}_{\BMO_{\mu,2}}=\max\set{\Norm{f}_{\bmo_{\mu,2}},\Norm{f}_{\bmo_{\mu,\mathsf{d}}}}.$$
In addition, the John-Nirenberg inequality holds in $\BMO_{\mu,p}$ so that $\BMO_{\mu,p}=\BMO_{\mu,1}$ for $p\in[1,\infty)$ with equivalent norm; see e.g. \cite[Chapter 4]{Lo1993}. 
Moreover,  for all $f\in\BMO_{\mu,2}$ with norm 1, there exist positive constants $C,\delta_{0}$ such that and for each $Q\in\D$,
\begin{align}\label{John-Nirenberg inequality}
\mu\left(\set{x\in Q:|f(x)-f_{\wh{Q}}|>\alpha}\right)\leq Ce^{-\delta_{0}\alpha}\mu(Q),
\end{align}
and also
\begin{align*}
\mu\left(\set{x\in Q:|f(x)-f_{Q}|>\alpha}\right)\leq Ce^{-\delta_{0}\alpha}\mu(Q).
\end{align*}

Recall that the John–Nirenberg inequality is not valid in general for $\bmo_{\mu,p}$ in place
of $\BMO_{\mu,p}$. A counter-example is given in \cite[Example 2.17]{We1994}. This is why we could not replace in general the space $\BMO_{\mu,2}$ by $\bmo_{\mu,2}$.




\subsection{Symmetric function space}
It is well-known that a Orlicz function space is related to a Young function.
A function  $\Psi:[0,\infty)\to[0,\infty)$ is said to be a Young function if there exists a function $\phi:[0,\infty)\to[0,\infty)$ such that $$\Psi(t):=\int_{0}^{t}\phi(s)ds,\quad t\geq0,$$
where $\phi$ is right-continuous and non-decreasing satisfying $\phi(0)=0$, $\displaystyle\lim_{s\to\infty}\phi(s)=\infty$ and $\phi(s)>0$ when $s>0$. Also say that $\Psi$ is a Young function generated by $\phi$. Moreover, we set $\displaystyle\psi(t)=\sup_{\phi(s)\leq t}s$ and denote $$\widetilde{\Psi}=\int_{0}^{t}\psi(s)ds,\quad t\geq0.$$
The function $\widetilde{\Psi}$ is called the complementary Young function of $\Psi$.
Given a pair  $(\Psi,\widetilde{\Psi})$ of Young function and its complementary, it has below property, for all $t>0$,
\begin{align*}
t\leq\Psi(t)\rp\widetilde{\Psi}(t)\rp\leq 2t.
\end{align*}  
 
Letting $\Psi:[0,\infty)\to[0,\infty)$ be a Young function, the $\Psi$-average of function $f$ over a cube $Q$ is defined by Luxemburg norm $$\Norm{f}_{L^{\Psi}(Q,d\mu)}=\inf\set{t>0:\frac{1}{\mu(Q)}\int_{Q}\Psi(\frac{|f(y)|}{t})d\mu(y)\leq1}.$$
Indeed, this is a local form of the Luxemburg norm. Recalling the equivalence for the Luxemburg norm (see \cite[Remarks 4.8.9]{PKJF-book}), one has a useful equivalent norm in calculations, i.e. there is a positive constant $C$ such that
\begin{align}\label{equivalent norm}
\Norm{f}_{L^{\Psi}(Q,d\mu)}
\leq\inf_{t>0}\set{t+\frac{t}{\mu(Q)}\int_{Q}\Psi(\frac{|f(y)|}{t})d\mu(y)}
\leq C\Norm{f}_{L^{\Psi}(Q,d\mu)}.
\end{align} 
In addition, we have a generalized H\"{o}lder's inequality.
\begin{lemma}\label{Holder inequality}
Let $\mu$ be a complete $\sigma$-finite measure and let $\widetilde{\Psi}$ be the complementary Young function of $\Psi$. For $Q\subset\RR$, if $f\in L^{\Psi}(Q,d\mu)$ and $g\in L^{\widetilde{\Psi}}(Q,d\mu)$, we have
\begin{align*}
\frac{1}{\mu(Q)}\int_{Q}|f(y)g(y)|d\mu(y)
\leq\Norm{f}_{L^{\Psi}(Q,d\mu)}\Norm{g}_{L^{\widetilde{\Psi}}(Q,d\mu)}.
\end{align*}
\end{lemma}
\begin{proof}
This is a direct result in combination of \cite[Theorem 4.7.5]{PKJF-book} and  \cite[Theorem 4.8.5]{PKJF-book}.
\end{proof}
In the proof of $L\log L$ estimate of the commutator, we pay our attention to the Young function $\Phi(t)=t(1+\log^{+}t)$, $t>0$. Its complementary Young function is given by $\widetilde{\Phi}(t)\simeq e^{t}$; see e.g.  \cite[Section 4.3]{PKJF-book}.



\section{$L_{p}$ estimate}\label{sec-Lpe}
In this section, we firstly prove \Cref{Lp estimate} with $p=2$ which will help us to show \Cref{endpoint estimate}. Assuming \Cref{endpoint estimate}, we then deduce \Cref{Lp estimate} for any $1<p<2$.


\subsection{Proof of $L_{2}$-boundedness}
Let us firstly prove the most important case the $L_{2}$ boundedness of $[M_{\varphi},H_{\D}]$ of \Cref{Lp estimate}.


\begin{lemma}\label{L2 estimate paraproduct}
Let $\mu$ be a positive Borel measure on $\RR$ with $0<\mu(Q)<\infty$ for $Q\in\D$ and $\varphi\in\bmo_{\mu,2}$. Define
\begin{align*}
\pi_{\varphi}f(x)
=\sum_{Q\in\D}f_{Q}\inner{\varphi,h_{Q}}h_{Q}(x),\quad f\in L_{2}(\mu).
\end{align*}
The operator $\pi_{\varphi}$ is bounded on $L_{2}(\mu)$ with upper bound $2\Norm{\varphi}_{\bmo_{\mu,2}}$.
\end{lemma}
\begin{proof}
By \cite[Lemma 5.10]{LMP2014}, we have
\begin{align*}
\Norm{\pi_{\varphi}f}_{L_{2}(\mu)}
\leq2\Norm{\set{\inner{\varphi,h_{Q}}}_{Q\in\D}}_{Carl(\mu)}\Norm{f}_{L_{2}(\mu)},
\end{align*}
where 
\begin{align*}
\Norm{\set{\inner{\varphi,h_{Q}}}_{Q\in\D}}_{Carl(\mu)}:=
\sup_{Q\in\D,\mu(Q)>0}\left(\frac{\sum_{Q'\in\D(Q)}|\inner{\varphi,h_{Q'}}|^{2}}{\mu(Q)}\right)^{\frac{1}{2}}.
\end{align*}
Since $\set{h_{Q}}_{Q\in\D}$ is an orthonormal bases and each $h_{Q}$ is cancelled over $Q$, it follows that
\begin{align*}
\sum_{Q'\in\D(Q)}|\inner{\varphi,h_{Q'}}|^{2}
&=\sum_{Q'\in\D(Q)}|\inner{(\varphi-\varphi_{Q})\chi_{Q},h_{Q'}}|^{2}\\
&\leq\Norm{(\varphi-\varphi_{Q})\chi_{Q}}_{L_{2}(\mu)}^{2}\leq\Norm{\varphi}_{\bmo_{\mu,2}}^{2}\mu(Q).
\end{align*}
Therefore, $$\Norm{\pi_{\varphi}f}_{L_{2}(\mu)}
\leq2\Norm{\varphi}_{\bmo_{\mu,2}}\Norm{f}_{L_{2}(\mu)}.$$
This is the desired assertion.

\end{proof}


\begin{remark}\label{remark on paraproduct}
\begin{enumerate}[\rm(1)]
\item Note that the operator $\pi_{\varphi}$ is bounded from $L_{1}(\mu)$ to $L_{1,\infty}(\mu)$ controlled by $\Norm{\set{\inner{\varphi,h_{Q}}}_{Q\in\D}}_{Carl(\mu)}$ (see \cite[Theorem 5.8 (i)]{LMP2014}).
Standard interpolation and duality argument gives us that  $\pi_{\varphi}$ is bounded on $L_{p}(\mu)$, $1<p\leq2$; $\pi_{\varphi}^{*}$ is bounded on $L_{q}(\mu)$, $2\leq q<\infty$. 

\item However, $\pi_{\varphi}$ is not bounded on $L_{q}(\mu)$ for $2< q<\infty$ and $\pi_{\varphi}^{*}$ is not bounded on $L_{p}(\mu)$ for $1<p<2$ except a very good measure $\mu$, e.g. dyadically doubling measure; see \cite[Theorem 5.8 (iii)]{LMP2014}.
\end{enumerate}
\end{remark}


For later convenience, we denote $\DD_{Q}f=\inner{f,h_{Q}}h_{Q}$ for $Q\in\D$ and $f\in L_{2}(\mu)$. Such an operator $\DD_{Q}$ is a projection on $L_{2}(\mu)$ and it admits the decomposition
\begin{align*}
\DD_{Q}f=\EE_{k}f\cdot\chi_{Q}-\EE_{k-1}f\cdot\chi_{Q}
=\EE_{Q}f-\sum_{Q'\in\D_{1}(Q)}\EE_{Q'}f
,\quad\mbox{if}\quad Q\in\D_{k},
\end{align*}
where $\EE_{Q}f:=\EE_{k}f\cdot\chi_{Q}$ if $Q\in\D_{k}$ and $\EE_{k}$ is given as in \eqref{expectation definition}. This in fact is the martingale difference defined on $Q$; for more details, we refer to \cite[Section 5.3]{LMP2014}.


\begin{lemma}\label{L2 estimate ad-paraproduct}
Let $\mu$ is a positive Borel measure on $\RR$ with $0<\mu(Q)<\infty$ for $Q\in\D$ and $\varphi\in\bmo_{\mu,\mathsf{d}}$. Define
\begin{align*}
\Lambda^{0}_{\varphi}f=\sum_{Q\in\D}\DD_{Q}\left(\DD_{Q}\varphi\cdot\DD_{Q}f\right),
\quad f\in L_{2}(\mu).
\end{align*}
The operator $\Lambda^{0}_{\varphi}$ is bounded on $L_{2}(\mu)$ with upper bound $\Norm{\varphi}_{\bmo_{\mu,\mathsf{d}}}$.
\end{lemma}
\begin{proof}
For each $Q\in\D$ and $f,g\in L_{2}(\mu)$, obviously $\inner{\DD_{Q}f,g}=\inner{f,\DD_{Q}g}$ since $\inner{\EE_{Q}f,g}=\inner{f,\EE_{Q}g}$. Therefore, 
\begin{align*}
\inner{\Lambda^{0}_{\varphi}f,g}
=\sum_{Q\in\D}\inner{\DD_{Q}\left(\DD_{Q}\varphi\cdot\DD_{Q}f\right),g}
&=\sum_{Q\in\D}\inner{\DD_{Q}\varphi\cdot\DD_{Q}f,\DD_{Q}g}\\
&\leq\left(\sum_{Q\in\D}\Norm{\DD_{Q}\varphi}_{L_{\infty}(\mu)}^{2}\Norm{\DD_{Q}f}_{L_{2}(\mu)}^{2}\right)^{\frac{1}{2}}
\left(\sum_{Q\in\D}\Norm{\DD_{Q}g}_{L_{2}(\mu)}^{2}\right)^{\frac{1}{2}}.
\end{align*}
Observing that
\begin{align*}
\Norm{\DD_{Q}\varphi}_{L_{\infty}(\mu)}
\leq\sup_{k\in\ZZ}\Norm{\EE_{k}\varphi-\EE_{k-1}\varphi}_{L_{\infty}(\mu)}
=\Norm{\varphi}_{\bmo_{\mu,\infty}}
\end{align*}
and (this is due to the orthogonality of martingale differences)
\begin{align*}
\left(\sum_{Q\in\D}\Norm{\DD_{Q}g}_{L_{2}(\mu)}^{2}\right)^{\frac{1}{2}}=\Norm{g}_{L_{2}(\mu)},\quad
\left(\sum_{Q\in\D}\Norm{\DD_{Q}f}_{L_{2}(\mu)}^{2}\right)^{\frac{1}{2}}=\Norm{f}_{L_{2}(\mu)},
\end{align*}
we arrive at
\begin{align*}
|\inner{\Lambda^{0}_{\varphi}f,g}|
\leq\Norm{\varphi}_{\bmo_{\mu,\mathsf{d}}}\Norm{f}_{L_{2}(\mu)}\Norm{g}_{L_{2}(\mu)}.
\end{align*}
Since $0<\mu(Q)<\infty$ for $Q\in\D$, this implies that $\mu$ is a $\sigma$-finite measure and therefore the dual representation holds for $L_{p}(\mu)$ with $1\leq p<\infty$.
Finally, the desired assertion follows from the dual representation of $L_{2}(\mu)$.
\end{proof}





%This proof is similar to \cite[Theorem A]{BAPW2024} but our Hilbert transform and measure condition are different. So it is valuable to write it down here for the completeness.

\begin{proposition}\label{L2 estimate}
Let $\mu$ be a $m$-increasing Borel measure on $\RR$ with $0<\mu(Q)<\infty$ for each $Q\in\D$.  If $\varphi\in\BMO_{\mu,2}$, then 
\begin{align*}
\Norm{[M_{\varphi},H_{\D}]}_{L_{2}(\mu)\rightarrow L_{2}(\mu)}
\leq C\Norm{\varphi}_{\BMO_{\mu,2}},
\end{align*}
for some positive constant $C$ independent of $\varphi$.
\end{proposition}
\begin{proof}
Take $f\in L_2(\mu)$. We write $f$ as the sum of martingale differences:
$$f=\sum_{k\in\mathbb{Z}}\mathbb{E}_k(f)-\mathbb{E}_{k-1}(f)=\sum_{Q\in\D}\DD_{Q}f.$$
Similarly, write
 $$\displaystyle\varphi=\sum_{Q\in\D}\DD_{Q}\varphi.$$
Both convergence can be understood in the sense of $L_{2}(\mu)$-norm.
Therefore, we apply \cite[Proposition 1.5]{Tr2013} to get
\begin{align*}
M_{\varphi}f=\pi_{\varphi}f+\Lambda^{0}_{\varphi}f+\Lambda^{1}_{\varphi}f+\pi_{\varphi}^{*}f,
\end{align*}
where
%\begin{align*}
%\pi_{\varphi}^{(*)}f(x)=\sum_{Q\in\D}\inner{\varphi,h_{Q}}\inner{f,h_{Q}}h_{Q}(x)^{2},
%\end{align*}
\begin{align*}
\pi_{\varphi}f
=\sum_{Q\in\D}\DD_{Q}\varphi\cdot \EE_{Q}f,
%=\sum_{Q\in\D}\inner{\varphi,h_{Q}}h_{Q}(x)\left(\sum_{J\in\D:J\supsetneq Q}\inner{f,h_{J}}h_{J}(x)\right)
%=\sum_{Q\in\D}f_{Q}\inner{\varphi,h_{Q}}h_{Q}(x)
\end{align*}
and $\pi_{\varphi}^{*}$ is the $L_{2}(\mu)$-adjoint of $\pi_{\varphi}$, and
\begin{align*}
\Lambda^{0}_{\varphi}f
=\sum_{Q\in\D}\DD_{Q}\left(\DD_{Q}\varphi\cdot\DD_{Q}f\right),\quad
\Lambda^{1}_{\varphi}f
=\sum_{Q\in\D}\EE_{Q}\varphi\cdot\DD_{Q}f.
%=\sum_{Q\in\D}\inner{f,h_{Q}}h_{Q}(x)\left(\sum_{J\in\D:J\supsetneq Q}\inner{\varphi,h_{J}}h_{J}(x)\right)
%=\sum_{Q\in\D}\varphi_{Q}\inner{f,h_{Q}}h_{Q}(x).
\end{align*}
The operators $\pi_{\varphi}$ is called paraproduct in \cite{Tr2013}. 
Since $\pi_{\varphi}^{*}$ is the $L_{2}(\mu)$-adjoint of $\pi_{\varphi}$, by \Cref{L2 estimate paraproduct}, $\pi_{\varphi}$ and $\pi_{\varphi}^{*}$ are bounded on $L_{2}(\mu)$ with
\begin{align*}
\Norm{\pi_{\varphi}}_{L_{2}(\mu)\rightarrow L_{2}(\mu)}
\lesssim\Norm{\varphi}_{\bmo_{\mu,2}}
\leq\Norm{\varphi}_{\BMO_{\mu}}
\quad\mbox{and}\quad
\Norm{\pi_{\varphi}^{*}}_{L_{2}(\mu)\rightarrow L_{2}(\mu)}
\lesssim\Norm{\varphi}_{\bmo_{\mu,2}}
\leq\Norm{\varphi}_{\BMO_{\mu}}.
\end{align*}
By \Cref{L2 estimate ad-paraproduct}, $\Lambda^{0}_{\varphi}$ is also bounded on $L_{2}(\mu)$ with
\begin{align*}
\Norm{\Lambda^{0}_{\varphi}}_{L_{2}(\mu)\rightarrow L_{2}(\mu)}
\leq\Norm{\varphi}_{\bmo_{\mu,\mathsf{d}}}
\leq \Norm{\varphi}_{\BMO_{\mu}}.
\end{align*}
Note that, in \cite[Theorem 2.5]{LMP2014}, the operator $H_{\D}$ is bounded on $L_{p}(\mu)$ for $1<p\leq2$. Therefore,
\begin{align*}
\Norm{[\pi_{\varphi},H_{\D}]}_{L_{2}(\mu)\rightarrow L_{2}(\mu)}
\lesssim \Norm{\varphi}_{\BMO_{\mu,2}},\quad
\Norm{[\pi_{\varphi}^{*},H_{\D}]}_{L_{2}(\mu)\rightarrow L_{2}(\mu)}
\lesssim \Norm{\varphi}_{\BMO_{\mu,2}}
\end{align*} 
and
\begin{align*}
\Norm{[\Lambda^{0}_{\varphi},H_{\D}]}_{L_{2}(\mu)\rightarrow L_{2}(\mu)}
\lesssim \Norm{\varphi}_{\BMO_{\mu,2}}.
\end{align*} 
So it is left to investigate the commutator $[\Lambda^{1}_{\varphi},H_{\D}]$.


By definition of $H_{\D}$, one has
\begin{align*}
H_{\D}(\Lambda^{1}_{\varphi}f)(x)=\sum_{Q\in\D}\sigma_{Q}\varphi_{\wh{Q}}\inner{f,h_{\wh{Q}}}h_{Q}(x)
\end{align*}
and
\begin{align*}
\Lambda^{1}_{\varphi}(H_{\D}f)(x)=\sum_{Q\in\D}\sigma_{Q}\varphi_{Q}\inner{f,h_{\wh{Q}}}h_{Q}(x).
\end{align*}
It gives that $$[\Lambda^{1}_{\varphi},H_{\D}]f(x)=
\sum_{Q\in\D}\sigma_{Q}\left(\varphi_{Q}-\varphi_{\wh{Q}}\right)\inner{f,h_{\wh{Q}}}h_{Q}(x).$$
Observing that $$\sup_{Q\in\D}|\sigma_{Q}\left(\varphi_{Q}-\varphi_{\wh{Q}}\right)|=\sup_k\sup_{Q\in \mathcal{D}_k}|\mathbb{E}_k(\varphi)\chi_Q-\mathbb{E}_{k-1}(\varphi)\chi_{Q}|=\|\varphi\|_{\bmo_{\mu,\mathsf{d}}}\leq\Norm{\varphi}_{\BMO_{\mu,2}},$$
we have
\begin{align*}
\Norm{[\Lambda^{1}_{\varphi},H_{\D}]f}_{L_{2}(\mu)}^{2}
&=\sum_{Q\in\D}|\sigma_{Q}\left(\varphi_{Q}-\varphi_{\wh{Q}}\right)|^{2}|\inner{f,h_{\wh{Q}}}|^{2}\\
&\leq\Norm{\varphi}_{\BMO_{\mu,2}}^{2}\sum_{Q\in\D}|\inner{f,h_{\wh{Q}}}|^{2}\\
&\leq2\Norm{\varphi}_{\BMO_{\mu,2}}^{2}\Norm{f}_{L_{2}(\mu)}^{2}.
\end{align*}
Therefore, the operator $[\Lambda^{1}_{\varphi},H_{\D}]$ is bounded on $L_{2}(\mu)$ with upper bound $\sqrt{2}\Norm{\varphi}_{\BMO_{\mu,2}}$. Finally, write the commutator $$[M_{\varphi},H_{\D}]=[\pi_{\varphi},H_{\D}]+[\Lambda^{0}_{\varphi},H_{\D}]
+[\Lambda^{1}_{\varphi},H_{\D}]+[\pi_{\varphi}^{*},H_{\D}].$$ The estimates on $L_{2}(\mu)$ of $[\pi_{\varphi},H_{\D}]$, $[\pi_{\varphi}^{*},H_{\D}]$, $[\Lambda^{0}_{\varphi},H_{\D}]$ and $[\Lambda^{1}_{\varphi},H_{\D}]$ announce that $[M_{\varphi},H_{\D}]$ is bounded on $L_{2}(\mu)$ with upper bound $\Norm{\varphi}_{\BMO_{\mu,2}}$ multiplying some universal constant.

\end{proof}

\begin{remark}
In viewing the proof of \Cref{L2 estimate}, the $L_{p}(\mu)$-boundedness $(1<p<2)$ of $[M_{\varphi},H_{\D}]$ is more difficult since $\pi_{\varphi}^{*}$ is not bounded on $L_{p}(\mu)$ with $1<p<2$ when $\mu$ is a non-doubling measure.
\end{remark}

\subsection{Proof of $L_{p}$-boundedness}
In this subsection, we prove \Cref{Lp estimate} under the assumption that we already have \Cref{endpoint estimate}. To prove $L_{p}$-boundedness $(1<p<2)$, we need a delicate estimate in calculus.

\begin{lemma}\label{integral estimate lemma}
Let $A>0$ and let $p>1$. There is a constant $C_{p}>0$ only depending on $p$ satisfying that
\begin{align*}
\int_{0}^{A}\log(1+\frac{A}{\lambda})\lambda^{p-2}d\lambda\leq C_{p}A^{p-1}.
\end{align*}
\end{lemma}
\begin{proof}
Letting $\lambda=\frac{A}{t}$ and integrating by parts, we have
\begin{align}\label{integral estimate}
\nonumber\int_{0}^{A}\log(1+\frac{A}{\lambda})\lambda^{p-2}d\lambda
&=A^{p-1}\int_{1}^{\infty}t^{-p}\log(1+t)dt\\
&=\frac{\log(2)}{p-1}A^{p-1}+
\frac{A^{p-1}}{p-1}\int_{1}^{\infty}t^{-p+1}\frac{1}{1+t}dt.
\end{align}
For the second factor in the last line on the right hand side, let us decompose the interval $(1,\infty)$ into pieces as follows
\begin{align*}
\int_{1}^{\infty}t^{-p+1}\frac{1}{1+t}dt
&=\sum_{k=0}^{\infty}\int_{2^k}^{2^{k+1}}t^{-p+1}\frac{1}{1+t}dt\\
&\leq \sum_{k=0}^{\infty}\int_{2^{k}}^{2^{k+1}}(2^{k})^{-p+1}\frac{1}{1+2^{k}}dt\leq\sum_{k=0}^{\infty}2^{k(1-p)}=:c_{p},
\end{align*}
where the series is convergent since $p>1$. Together this estimate with \eqref{integral estimate}, we obtain the desired assertion.

\end{proof}


To prove \Cref{Lp estimate} for $1<p<2$, let us firstly admit \Cref{endpoint estimate} whose proof will be provided in \Cref{sec-end}. 

\begin{proof}[Proof of \Cref{Lp estimate}]
 It suffices to consider the cases for $1<p<2$ since the case $p=2$ is included in \Cref{L2 estimate}. Without loss of generality, assume $\Norm{\varphi}_{\BMO_{\mu,2}}=1$ firstly. For $f\in L_{p}(\mu)$ and any $\lambda>0$, decompose $f=f_{\lambda}+f^{\lambda}$, where
\begin{equation*}
f_{\lambda}(x)=\left\{
\begin{aligned}
f(x),\quad&|f(x)|>\lambda\\
0,\quad&|f(x)|\leq\lambda;
\end{aligned}
\right.
\quad\mbox{and}\quad
f^{\lambda}(x)=\left\{
\begin{aligned}
0,\quad&|f(x)|>\lambda\\
f(x),\quad&|f(x)|\leq\lambda.
\end{aligned}
\right.
\end{equation*}
By distribution formula \cite[Proposition 1.1.4]{Gr-book1}, one has
\begin{align*}
\Norm{[M_{\varphi},H_{\D}]f}_{L_{p}(\mu)}^{p}
&=p2^{p-1}\int_{0}^{\infty}\mu\left(\set{x\in\RR:|[M_{\varphi},H_{\D}]f(x)|>2\lambda}\right)
\lambda^{p-1}d\lambda\\
&\lesssim\int_{0}^{\infty}\mu\left(\set{x\in\RR:|[M_{\varphi},H_{\D}]f^{\lambda}(x)|>\lambda}\right)
\lambda^{p-1}d\lambda\\
&\quad+\int_{0}^{\infty}\mu\left(\set{x\in\RR:|[M_{\varphi},H_{\D}]f_{\lambda}(x)|>\lambda}\right)
\lambda^{p-1}d\lambda\\
&\lesssim\int_{0}^{\infty}\Norm{f^{\lambda}}_{L_{2}(\mu)}^{2}\lambda^{p-3}d\lambda
+\int_{0}^{\infty}\int_{\RR}\frac{|f_{\lambda}(x)|}{\lambda}\left(1+\log^{+}(\frac{|f_{\lambda}(x)|}{\lambda})\right)
d\mu(x)\lambda^{p-1}d\lambda,
\\&=\int_{0}^{\infty}\Norm{f^{\lambda}}_{L_{2}(\mu)}^{2}\lambda^{p-3}d\lambda+\mathrm{II}
\end{align*}
where we used \Cref{L2 estimate} and \Cref{endpoint estimate} in the last inequality.
Note that, for $1<p<2$,
\begin{align*}
\int_{0}^{\infty}\Norm{f^{\lambda}}_{L_{2}(\mu)}^{2}\lambda^{p-3}d\lambda
=\int_{\RR}\int_{|f(x)|}^{\infty}\lambda^{p-3}d\lambda\cdot|f(x)|^{2}d\mu(x)
=\frac{1}{2-p}\Norm{f}_{L_{p}(\mu)}^{p}.
\end{align*}
On the other hand side, applying \Cref{integral estimate lemma} to $A=f(x)$ with $f(x)\neq0$, it follows that
\begin{align*}
\mathrm{II}&= \int_{\RR}\int_{0}^{|f(x)|}
\lambda^{p-2}d\lambda \cdot |f(x)|d \mu(x)+ \int_{\RR}\int_{0}^{|f(x)|}\left(1+\log^{+}(\frac{|f_{\lambda}(x)|}{\lambda})\right)
\lambda^{p-2}d\lambda d \cdot |f_{\lambda}(x)|\mu(x)\\
&\leq\frac{1}{p-1}\Norm{f}_{L_{p}(\mu)}^{p}
+\int_{\RR}\int_{0}^{|f(x)|}\lambda^{p-2}\log\left(1+\frac{|f(x)|}{\lambda}\right)d\lambda\cdot|f(x)|d\mu(x)\\
&\lesssim\Norm{f}_{L_{p}(\mu)}^{p}.
\end{align*}
Therefore, for $1<p<2$ and $\Norm{\varphi}_{\BMO_{\mu,2}}=1$, we have
\begin{align*}
\Norm{[M_{\varphi},H_{\D}]f}_{L_{p}(\mu)}\lesssim\Norm{f}_{L_{p}(\mu)},
\end{align*}
where the constant only depends  on $p$. Now, when $\Norm{\varphi}_{\BMO_{\mu,2}}\neq1$, let us write $$[M_{\varphi},H_{\D}]f=[M_{\frac{\varphi}{\Norm{\varphi}_{\BMO_{\mu,2}}}},H_{\D}](\Norm{\varphi}_{\BMO_{\mu,2}}f),$$
which gives the assertion of $[M_{\varphi},H_{\D}]$ for $1<p<2$ from the previous argument. The assertion of $[M_{\varphi},H_{\D}^{*}]$ immediately follows from standard  duality argument. The proof is complete.
\end{proof}


\section{Endpoint estimate}\label{sec-end}
In this section, we show \Cref{endpoint estimate} by using the Calder\'{o}n-Zygmund decomposition and  \Cref{L2 estimate}. In \Cref{sec-nece lemmas}, we present useful lemmas, and the detailed proof of \Cref{endpoint estimate} is given in \Cref{sec-detailed proof}.

\subsection{Necessary lemmas}\label{sec-nece lemmas}
The following two   lemmas will be used later.


\begin{lemma}\label{estimate to T32}
Let $\varphi$ and $f$ be as in \Cref{endpoint estimate}, and let $\Omega_{\alpha}=\cup_jQ_j$ be as in \Cref{CZ decomposition} at level $\alpha=\frac{\lambda}{\Norm{\varphi}_{\BMO_{\mu,2}}}$, $\lambda>0$. Denote $\mathcal{Q}=\{Q_j\}_j$. For each $Q\in \mathcal{Q}$, we have
\begin{align*}
\frac{1}{\lambda}\int_{Q}\left|(\varphi(y)-\varphi_{Q})f(y)\right|d\mu(y)
\leq2\int_{Q}\frac{\Norm{\varphi}_{\BMO_{\mu,2}}|f(y)|}{\lambda}
\left(1+\log^{+}(\frac{\Norm{\varphi}_{\BMO_{\mu,2}}|f(y)|}{\lambda})\right)d\mu(y).
\end{align*}
\end{lemma}
\begin{proof}
Recall that the complementary Young function of $\Phi(t)=t(1+\log^{+}t)$ is given by $\widetilde{\Phi}(t)\simeq e^{t}$ (see \cite[Section 4.3]{PKJF-book}). By generalized H\"{o}lder's inequality in \Cref{Holder inequality}, we have
\begin{align*}
\int_{Q}\left|(\varphi(y)-\varphi_{Q})f(y)\right|d\mu(y)
&=\frac{1}{\Norm{\varphi}_{\BMO_{\mu,2}}}
\int_{Q}|\varphi(y)-\varphi_{Q}|\cdot\Norm{\varphi}_{\BMO_{\mu,2}}|f(y)|d\mu(y)\\
&\leq\frac{1}{\Norm{\varphi}_{\BMO_{\mu,2}}}
\Norm{\varphi-\varphi_{Q}}_{L^{\exp}(Q,\mu)}
\Norm{\Norm{\varphi}_{\BMO_{\mu,2}}f}_{L^{\Phi}(Q,\mu)}\mu(Q).
\end{align*}
By the John-Nirenberg inequality \eqref{John-Nirenberg inequality}, one has
\begin{align}\label{exp J-N inequality}
\Norm{\varphi-\varphi_{Q}}_{L^{\exp}(Q,\mu)}\leq C\Norm{\varphi}_{\BMO_{\mu,2}}.
\end{align}
Indeed, for an undetermined positive constant $C$, by John-Nirenberg inequality,
\begin{align*}
\frac{1}{\mu(Q)}\int_{Q}e^{\frac{|\varphi(y)-\varphi_{Q}|}{C\Norm{\varphi}_{\BMO_{\mu,2}}}}d\mu(y)
&=\frac{1}{\mu(Q)}\int_{0}^{\infty}e^{t}
\mu\left(\{y\in Q:|\varphi(y)-\varphi_{Q}|\geq C\Norm{\varphi}_{\BMO_{\mu,2}}t\}\right)dt\\
&\leq\int_{0}^{\infty}e^{(1-\delta C)t}dt.
\end{align*}
Now, we can pick a large positive constant $C$ such that $\delta C>1$ and the last integral is less that 1.  This implies \eqref{exp J-N inequality} holds.
For $\Phi(t)=t(1+\log^{+}t)$, recalling \eqref{equivalent norm}, we have
\begin{align*}
\frac{1}{\lambda}\Norm{\Norm{\varphi}_{\BMO_{\mu,2}}f}_{L^{t(1+\log^{+} t)}(Q,\mu)}\mu(Q)
&\leq\frac{1}{\lambda}\mu(Q)\inf_{t>0}
\set{t+\frac{t}{\mu(Q)}\int_{Q}\Phi(\frac{\Norm{\varphi}_{\BMO_{\mu,2}}|f(y)|}{t})d\mu(y)}\\
&\leq \mu(Q)+\int_{Q}\Phi(\frac{\Norm{\varphi}_{\BMO_{\mu,2}}|f(y)|}{\lambda})d\mu(y)\\
&\leq \frac{\Norm{\varphi}_{\BMO_{\mu,2}}}{\lambda}\int_{Q}|f(y)|d\mu(y)
+\int_{Q}\Phi(\frac{\Norm{\varphi}_{\BMO_{\mu,2}}|f(y)|}{\lambda})d\mu(y)\\
&\leq2\int_{Q}\Phi(\frac{\Norm{\varphi}_{\BMO_{\mu,2}}|f(y)|}{\lambda})d\mu(y),
\end{align*}
where in the last second inequality we use the fact from Calder\'{o}n-Zygmund decomposition in \Cref{CZ decomposition} that $$\frac{\Norm{\varphi}_{\BMO_{\mu,2}}}{\mu(Q)}\int_{Q}|f(y)|d\mu(y)>\lambda.$$
It gives that
\begin{align*}
\frac{1}{\lambda}\Norm{\Norm{\varphi}_{\BMO_{\mu,2}}f}_{L\log L(\mu),Q}|Q|
\leq2\int_{Q}\frac{\Norm{\varphi}_{\BMO_{\mu,2}}|f(y)|}{\lambda}\left(1+\log^{+}
(\frac{\Norm{\varphi}_{\BMO_{\mu,2}}|f(y)|}{\lambda})\right)d\mu(y).
\end{align*}
This inequality and estimate \eqref{exp J-N inequality} imply the assertion.

\end{proof}



\begin{lemma}\label{estimate to T41}
Let $\varphi$ and $f$ be as in \Cref{endpoint estimate}. Let $\set{\beta_{j}}_{j}$ and $\set{Q_{j}}_{j}\subset\D$ be as  in \Cref{CZ decomposition} at level $\alpha=\frac{\lambda}{\Norm{\varphi}_{\BMO_{\mu,2}}}$,  $0<\lambda<\infty$. We have
\begin{align*}
\mathsf{B}_{\lambda}:=\mu\left(\set{y\in\RR\setminus\Omega_{\alpha}:
\left|\sum_{j}(\varphi(y)-\varphi_{\wh{Q}_{j}})H_{\D}\beta_{j}(y)\right|>\frac{\lambda}{6}}\right)
\lesssim\frac{\Norm{\varphi}_{\BMO_{\mu,2}}}{\lambda}\Norm{f}_{L_{1}(\mu)}.
\end{align*}
\end{lemma}
\begin{proof}
For each $j$, $\beta_{j}$ has vanishing integrals with support $\wh{Q}_{j}$, is a constant on each dyadic children of $\wh{Q}_{j}$, Then, we have
$$\inner{\beta_{j},h_{Q}}=
\sigma(Q_{j})(f_{Q_{j}}-f_{\wh{Q}_{j}})m(\wh{Q}_{j})^{\frac{1}{2}}\delta_{\wh{Q}_{j},Q},$$
which further implies
$$H_{\D}\beta_{j}(y)=\sum_{Q\in\D}\sigma(Q)\inner{\beta_{j},h_{\wh{Q}}}h_{Q}(y)
=(f_{Q_{j}}-f_{\wh{Q}_{j}})m(\wh{Q}_{j})^{\frac{1}{2}}(h_{Q_{j}}(y)-h_{Q_{j}^{b}}(y)),$$
where $Q_{j}^{b}=\wh{Q}_{j}\setminus Q_{j}\in\D$ is the dyadic brother of $Q\in\D$. 
It follows that each $H_{\D}\beta_{j}$ is supported in $\wh{Q}_{j}$, and then
\begin{equation}\label{estimate T41}
\begin{split}
\mathsf{B}_{\lambda}
&\lesssim\frac{1}{\lambda}
\sum_{j}\int_{\wh{Q}_{j}}|(\varphi(y)-\varphi_{\wh{Q}_{j}})H_{\D}\beta_{j}(y)|d\mu(y)\\
&\lesssim
\frac{1}{\lambda}\sum_{j}|f_{Q_{j}}-f_{\wh{Q}_{j}}|m(\wh{Q}_{j})^{\frac{1}{2}}\cdot
\int_{\wh{Q}_{j}}|\varphi(y)-\varphi_{\wh{Q}_{j}}|\cdot|h_{Q_{j}}(y)-h_{Q_{j}^{b}}(y)|d\mu(y).
\end{split}
\end{equation}
Next, let us investigate the integrand on the last line. For $Q_{j}\in\D_{j}$, we know its father $\wh{Q}_{j}\in\D_{j-1}$ and  $h_{Q_{j}},h_{Q_{j}^{b}}\in\F_{j+1}$ . Thus,
\begin{align*}
\int_{\wh{Q}_{j}}|\varphi(y)-\varphi_{\wh{Q}_{j}}|\cdot|h_{Q_{j}}(y)-h_{Q_{j}^{b}}(y)|d\mu(y)
&=\int_{\wh{Q}_{j}}|\varphi(y)-\EE_{j-1}(\varphi)(y)|\cdot|h_{Q_{j}}(y)-h_{Q_{j}^{b}}(y)|d\mu(y)\\
&\leq\int_{\wh{Q}_{j}}|\varphi(y)-\EE_{j+1}(\varphi)(y)||h_{Q_{j}}(y)-h_{Q_{j}^{b}}(y)|d\mu(y)\\
&\quad+\int_{\wh{Q}_{j}}|\EE_{j+1}(\varphi)(y)-\EE_{j}(\varphi)(y)||h_{Q_{j}}(y)-h_{Q_{j}^{b}}(y)|d\mu(y)\\
&\quad+\int_{\wh{Q}_{j}}|\EE_{j}(\varphi)(y)-\EE_{j-1}(\varphi)(y)||h_{Q_{j}}(y)-h_{Q_{j}^{b}}(y)|d\mu(y)\\
&=\mathrm{I}_1+\mathrm{II}_2+\mathrm{III}_3.
\end{align*}
{\color{red}I stop here.}
We next estimate the three terms on the right hand side separately.
Observing that $Q_{j},Q_{j}^{b}\subset \wh{Q}_{j}$ and measure $\mu$ is $m$-increasing, it follows from \eqref{hQ norm} that 
\begin{align}\label{hQ-hQ integral}
\int_{\wh{Q}_{j}}|h_{Q_{j}}-h_{Q_{j}^{b}}|d\mu(y)
\leq2m(Q_{j})^{\frac{1}{2}}+2m(Q_{j}^{b})^{\frac{1}{2}}
\lesssim m(\wh{Q}_{j})^{\frac{1}{2}}.
\end{align}
We have
\begin{align*}
\int_{\wh{Q}_{j}}|\EE_{j+1}(\varphi)(y)-\EE_{j}(\varphi)(y)||h_{Q_{j}}(y)-h_{Q_{j}^{b}}(y)|d\mu(y)
&\leq\Norm{\varphi}_{\bmo_{\mu,\infty}}\int_{\wh{Q}_{j}}|h_{Q_{j}}(y)-h_{Q_{j}^{b}}(y)|d\mu(y)\\
&\lesssim\Norm{\varphi}_{\bmo_{\mu,\infty}}m(\wh{Q}_{j})^{\frac{1}{2}}.
\end{align*}
Similarly,
\begin{align*}
\int_{\wh{Q}_{j}}|\EE_{j}(\varphi)(y)-\EE_{j-1}(\varphi)(y)||h_{Q_{j}}(y)-h_{Q_{j}^{b}}(y)|d\mu(y)
\lesssim\Norm{\varphi}_{\bmo_{\mu,\infty}}m(\wh{Q}_{j})^{\frac{1}{2}}.
\end{align*}
Since $h_{Q_{j}},h_{Q_{j}^{b}}\in\F_{j+1}$ and $\chi_{\wh{Q}_{j}}\in\F_{j-1}\subset\F_{j+1}$, it gives from \eqref{hQ-hQ integral} that
\begin{align*}
&\int_{\RR}|\varphi(y)-\EE_{j+1}(\varphi)(y)||h_{Q_{j}}(y)-h_{Q_{j}^{b}}(y)|\chi_{\wh{Q}_{j}}(y)d\mu(y)\\
&=\int_{\RR}\EE_{j+1}\left(|\varphi-\EE_{j+1}(\varphi)||h_{Q_{j}}-h_{Q_{j}^{b}}|\chi_{\wh{Q}_{j}}\right)(y)d\mu(y)\\
&=\int_{\RR}|h_{Q_{j}}(y)-h_{Q_{j}^{b}}(y)|\chi_{\wh{Q}_{j}}(y)\EE_{j+1}\left(|\varphi-\EE_{j+1}(\varphi)|\right)(y)d\mu(y)\\
&\leq\sup_{j}\Norm{\EE_{j+1}\left(|\varphi-\EE_{j+1}(\varphi)|\right)}_{\infty}\int_{\wh{Q}_{j}}|h_{Q_{j}}(y)-h_{Q_{j}^{b}}(y)|d\mu(y)\\
&\lesssim\Norm{\varphi}_{\BMO_{\mu,2}}m(\wh{Q}_{j})^{\frac{1}{2}}.
\end{align*}
In a sum, by \eqref{BMO intersection space},
\begin{align*}%\label{p-pQ hQ-hQ estimate}
\int_{\wh{Q}_{j}}|\varphi(y)-\varphi_{\wh{Q}_{j}}||h_{Q_{j}}(y)-h_{Q_{j}^{b}}(y)|d\mu(y)
\lesssim\Norm{\varphi}_{\BMO_{\mu,2}}m(\wh{Q}_{j})^{\frac{1}{2}}.
\end{align*}
Therefore, combining \eqref{estimate T41} with this estimate,
\begin{align*}
\mathsf{B}_{\lambda}
&\lesssim\frac{\Norm{\varphi}_{\BMO_{\mu,2}}}{\lambda}\sum_{j}|f_{Q_{j}}-f_{\wh{Q}_{j}}|m(\wh{Q}_{j})\\
&\lesssim\frac{\Norm{\varphi}_{\BMO_{\mu,2}}}{\lambda}\sum_{j}|f|_{Q_{j}}m(\wh{Q}_{j})\\
&\lesssim\frac{\Norm{\varphi}_{\BMO_{\mu,2}}}{\lambda}\sum_{j}\int_{Q_{j}}|f(y)|d\mu(y)\\
&\lesssim\frac{\Norm{\varphi}_{\BMO_{\mu,2}}}{\lambda}\Norm{f}_{L_{1}(\mu)}.
\end{align*}
where we use $|f|_{\wh{Q}_{j}}\leq\alpha<|f|_{Q_{j}}$ in the second line which can be discerned in \cite[(3.2) in page 57]{LMP2014} and we use $m(\wh{Q}_{j})\leq\mu(Q_{j})$ in the third line which can be found in \cite[page 51]{LMP2014}. 


\end{proof}




\subsection{Proof of \Cref{endpoint estimate}}\label{sec-detailed proof}

We are now ready to prove \Cref{endpoint estimate}.

\begin{proof}[Proof of \Cref{endpoint estimate}]
Fix $\lambda>0$. Decompose $\RR=\Omega_{\alpha}\bigcup(\RR\setminus\Omega_{\alpha})$ where $\displaystyle\Omega_{\alpha}=\bigcup_{j}Q_{j}$ is the set arising from the Calder\'{o}n-Zygmund decomposition for $f$ in \Cref{CZ decomposition} at level $\alpha=\frac{\lambda}{\Norm{\varphi}_{\BMO_{\mu,2}}}$.
Therefore, $f=g+b+\beta$ and
\begin{align*}
\mu\left(\set{y\in\RR:\left|[M_{\varphi},H_{\D}]f(y)\right|>\lambda}\right)
&\leq\mu\left(\set{y\in\Omega_{\alpha}:\left|[M_{\varphi},T]f(y)\right|>\lambda}\right)\\
&\quad+\mu\left(\set{y\in\RR\setminus\Omega_{\alpha}:\left|[M_{\varphi},H_{\D}]g(y)\right|>\frac{\lambda}{3}}\right)\\
&\quad+\mu\left(\set{y\in\RR\setminus\Omega_{\alpha}:\left|[M_{\varphi},H_{\D}]b(y)\right|>\frac{\lambda}{3}}\right)\\
&\quad+\mu\left(\set{y\in\RR\setminus\Omega_{\alpha}:\left|[M_{\varphi},H_{\D}]\beta(y)\right|>\frac{\lambda}{3}}\right)\\
&=:{\rm T}_{1}+{\rm T}_{2}+{\rm T}_{3}+{\rm T}_{4}.
\end{align*}


By \Cref{CZ decomposition} \eqref{Omega upper bd} and $\alpha=\frac{\lambda}{\Norm{\varphi}_{\BMO_{\mu,2}}}$, we have
\begin{align*}
{\rm T}_{1}\leq\mu(\Omega_{\alpha})
\lesssim\frac{\Norm{\varphi}_{\BMO_{\mu,2}}\Norm{f}_{L_{1}(\mu)}}{\lambda}.
\end{align*}



It follows from \Cref{L2 estimate} that $[M_{\varphi},H_{\D}]$ is bounded on $L_{2}(\mu)$ provided $\varphi\in\BMO_{\mu,2}$. Therefore, combining \Cref{CZ decomposition} \eqref{g upper bd},
\begin{align*}
{\rm T}_{2}
&\leq\frac{4}{\lambda^{2}}\int_{\RR\setminus\Omega_{\alpha}}|[M_{\varphi},H_{\D}]g(y)|^{2}d\mu(y)\\
&\lesssim\frac{\Norm{\varphi}_{\BMO_{\mu,2}}^{2}}{\lambda^{2}}\int_{\RR}|g(y)|^{2}d\mu(y)\lesssim\frac{\Norm{\varphi}_{\BMO_{\mu,2}}}{\lambda}\Norm{f}_{L_{1}(\mu)}.
\end{align*}


Now, let us estimate ${\rm T}_{3}$. Write 
$$b=\sum_{j}b_{j}=(f-f_{Q_{j}})\chi_{Q_{j}}.$$
Observing that each $b_{j}$ has vanishing integral and that $h_{Q}$ is constant on each $Q_{\pm}$, it is easy to see that $H_{\D}b_{j}(x)=0$ whenever $x\in \RR\setminus Q_{j}$. We have
\begin{align*}
[M_{\varphi},H_{\D}]b(y)&=\sum_{j}[M_{\varphi-\varphi_{Q_{j}}},H_{\D}]b_{j}(y)\\
&=\sum_{j}(\varphi(y)-\varphi_{Q_{j}})H_{\D}b_{j}(y)-\sum_{j}H_{\D}((\varphi-\varphi_{Q_{j}})b_{j})(y)\\
&=-\sum_{j}H_{\D}((\varphi-\varphi_{Q_{j}})b_{j})(y).
\end{align*}
Thus,
\begin{align*}
{\rm T}_{3}=
\mu\left(\set{y\in\RR\setminus\Omega_{\alpha}:\left|\sum_{j}H_{\D}((\varphi-\varphi_{Q_{j}})b_{j})(y)\right|>\frac{\lambda}{3}}\right).
\end{align*}
Since $H_{\D}$ is of type weak $(1,1)$ when $\mu$ is $m$-increasing (see \cite[Theorem 2.5]{LMP2014}), it follows that
\begin{align*}
{\rm T}_{3}
&\lesssim\frac{1}{\lambda}\int_{\RR}\left|\sum_{j}(\varphi(y)-\varphi_{Q_{j}})b_{j}(y)\right|d\mu(y)\\
&\lesssim\frac{1}{\lambda}\sum_{j}\int_{Q_{j}}\left|(\varphi(y)-\varphi_{Q_{j}})f_{Q_{j}}\right|d\mu(y)
+\frac{1}{\lambda}\sum_{j}\int_{Q_{j}}\left|(\varphi(y)-\varphi_{Q_{j}})f(y)\right|d\mu(y)\\
&=:{\rm T}_{31}+{\rm T}_{32}.
\end{align*}
By the non-overlapping property of $\{Q_{j}\}$, we have
\begin{align*}%\label{estimate iii21}
{\rm T}_{31}
&\leq\frac{1}{\lambda}\sum_{j}\frac{1}{\mu(Q_{j})}\int_{Q_{j}}|\varphi(y)-\varphi_{Q_{j}}|d\mu(y)
\int_{Q_{j}}|f(x)|d\mu(x)\\
&\leq\frac{\Norm{\varphi}_{\bmo_{\mu,2}}}{\lambda}\Norm{f}_{L_{1}(\mu)}\leq \frac{\Norm{\varphi}_{\BMO_{\mu,2}}}{\lambda}\Norm{f}_{L_{1}(\mu)}.
\end{align*}
As for ${\rm T}_{32}$, applying \Cref{estimate to T32} and using the fact that $\set{Q_{j}}_{j}$ are disjoint, one obtains
\begin{align*}%\label{estimate iii22}
{\rm T}_{32}&\leq2
\sum_{j}\int_{Q_{j}}\frac{\Norm{\varphi}_{\BMO_{\mu,2}}|f(y)|}{\lambda}
\left(1+\log^{+}(\frac{\Norm{\varphi}_{\BMO_{\mu,2}}|f(y)|}{\lambda})\right)d\mu(y)\\
&\leq2\int_{\RR}\frac{\Norm{\varphi}_{\BMO_{\mu,2}}|f(y)|}{\lambda}
\left(1+\log^{+}(\frac{\Norm{\varphi}_{\BMO_{\mu,2}}|f(y)|}{\lambda})\right)d\mu(y).
\end{align*}
Thus, 
\begin{align*}
{\rm T}_{3}\leq {\rm T}_{31}+{\rm T}_{32}\lesssim
\int_{\RR}\frac{\Norm{\varphi}_{\BMO_{\mu,2}}|f(y)|}{\lambda}\left(1+\log^{+}(\frac{\Norm{\varphi}_{\BMO_{\mu,2}}|f(y)|}{\lambda})\right)d\mu(y).
\end{align*}


Next, let us investigate ${\rm T}_{4}$. Recalling that $\displaystyle \beta=\sum_{j}\beta_{j}$, we write
$$[M_{\varphi},H_{\D}]\beta(y)=\sum_{j}[M_{\varphi-\varphi_{Q_{j}}},H_{\D}]\beta_{j}(y)
=\sum_{j}(\varphi(y)-\varphi_{\wh{Q}_{j}})H_{\D}\beta_{j}(y)-\sum_{j}H_{\D}((\varphi-\varphi_{\wh{Q}_{j}})\beta_{j})(y).$$
Thus,
\begin{align*}
{\rm T}_{4}
&\leq\mu\left(\set{y\in\RR\setminus\Omega_{\alpha}:\left|\sum_{j}(\varphi(y)-\varphi_{\wh{Q}_{j}})H_{\D}\beta_{j}(y)\right|>\frac{\lambda}{6}}\right)\\
&\quad+\mu\left(\set{y\in\RR\setminus\Omega_{\alpha}:\left|\sum_{j}H_{\D}((\varphi-\varphi_{\wh{Q}_{j}})\beta_{j})(y)\right|>\frac{\lambda}{6}}\right)\\
&=:{\rm T}_{41}+{\rm T}_{42}.
\end{align*}
By \Cref{estimate to T41}, we have
\begin{align*}
{\rm T}_{41}\lesssim\frac{\Norm{\varphi}_{\BMO_{\mu,2}}}{\lambda}\Norm{f}_{L_{1}(\mu)}.
\end{align*}
For ${\rm T}_{42}$, noting that $H_{\D}$ is of type weak $(1,1)$ when $\mu$ is $m$-increasing (see \cite[Theorem 2.5]{LMP2014}) and $\beta_{j}=(f_{Q_{j}}-f_{\wh{Q}_{j}})
\left(\chi_{Q_{j}}(x)-\frac{\mu(Q_{j})}{\mu(\wh{Q}_{j})}\chi_{\wh{Q}_{j}}(x)\right)$, it gives that
\begin{align*}
{\rm T}_{42}
&\lesssim\frac{1}{\lambda}\sum_{j}\int_{Q_{j}}|\varphi(y)-\varphi_{\wh{Q}_{j}}|\cdot
|f_{Q_{j}}-f_{\wh{Q}_{j}}|d\mu(y)\\
&\quad+\frac{1}{\lambda}\sum_{j}\int_{\wh{Q}_{j}}|\varphi(y)-\varphi_{\wh{Q}_{j}}|\cdot
|f_{Q_{j}}-f_{\wh{Q}_{j}}|\frac{\mu(Q_{j})}{\mu(\wh{Q}_{j})}d\mu(y).
\end{align*}
Recall that for each $j$, $|f|_{\wh{Q}_{j}}\leq\alpha<|f|_{Q_{j}}$; see \Cref{CZ decomposition}. Hence,
\begin{align*}
{\rm T}_{42}
&\lesssim\frac{2}{\lambda}\sum_{j}\int_{Q_{j}}|\varphi(y)-\varphi_{\wh{Q}_{j}}|\cdot
|f|_{Q_{j}}d\mu(y)
+\frac{2}{\lambda}\sum_{j}\int_{\wh{Q}_{j}}|\varphi(y)-\varphi_{\wh{Q}_{j}}|
\cdot|f|_{Q_{j}}\frac{\mu(Q_{j})}{\mu(\wh{Q}_{j})}d\mu(y)\\
&\lesssim\frac{1}{\lambda}\sum_{j}\frac{1}{\mu(Q_{j})}\int_{Q_{j}}|\varphi(y)-\varphi_{\wh{Q}_{j}}|
d\mu(y)\int_{Q_{j}}|f(x)|d\mu(x)\\
&\quad+\frac{1}{\lambda}\sum_{j}\frac{1}{\mu(\wh{Q}_{j})}\int_{\wh{Q}_{j}}|\varphi(y)-\varphi_{\wh{Q}_{j}}|d\mu(y)
\int_{Q_{j}}|f(x)|d\mu(x)\\
&\leq\frac{1}{\lambda}(\Norm{\varphi}_{\BMO_{\mu,2}}+\Norm{\varphi}_{\bmo_{\mu,2}})\Norm{f}_{L_{1}(\mu)}\leq \frac{2\Norm{\varphi}_{\BMO_{\mu,2}}}{\lambda} \Norm{f}_{L_{1}(\mu)},
\end{align*}
where the inequality is due to \eqref{BMO intersection space}. Combing the estimates on ${\rm T}_{41}$ and $\mathrm{T}_{42}$, we arrive at
\begin{align*}
{\rm T}_{4}\lesssim\frac{\Norm{\varphi}_{\BMO_{\mu,2}}}{\lambda}\Norm{f}_{L_{1}(\mu)}.
\end{align*} 



Finally, the desired assertion follows from the combination of estimates for ${\rm T}_{1}$, ${\rm T}_{2}$, ${\rm T}_{3}$ and ${\rm T}_{4}$.  
\end{proof}




\section{Compactness}
%\subsection{Definition }
%martingale version in dyadic filtration with Lebesgue measure: 
%$$\|f\|_{B_p}=\Big(\sum_{k\geq1}\sum_{I\in D_k} \Big[|I|^{-1}\|d_kf\|_p^p\Big] \Big)^{1/p}$$
%
%$d_nb= \sum_{I\in D_{n-1}} <b,h_I> h_I$

\subsection{Besov space}
For $0<p<\infty$, define the martingale Besov space by
\begin{align*}
\Norm{f}_{\dot{\htb}_{p}(\mu)}=
\left(\sum_{k\in\ZZ}\Norm{\beta_{k}^{\frac{1}{p}}\DD_{k}f}_{L_{p}(\mu)}^{p}\right)^{\frac{1}{p}},
\end{align*}
where $\displaystyle\beta_{k}=\sum_{Q\in\D_{k}}\mu(Q)\rp\chi_{Q}$. 
{\color{red}Prove such definition does not rely on choice of basis ?}


\begin{lemma}\label{besov equi-version}
For $0<p<\infty$,, we have
\begin{align*}
\Norm{f}_{\dot{\htb}_{p}(\mu)}=
\left(\sum_{Q\in\D}|\inner{f,h_{Q}}|^{p}m(Q)^{\frac{p}{2}}
\left(\frac{1}{\mu(Q_{-})^{p}}+\frac{1}{\mu(Q_{+})^{p}}\right)\right)^{\frac{1}{p}}.
\end{align*}
\end{lemma}
\begin{proof}
Recall that $\displaystyle\DD_{k}f=\sum_{Q\in\D_{k-1}}\inner{f,h_{Q}}h_{Q}$. By $h_{Q}=m(Q)^{\frac{1}{2}}\left(\frac{\chi_{Q_{-}}}{\mu(Q_{-})}-\frac{\chi_{Q_{+}}}{\mu(Q_{+})}\right)$ and definition of $\beta_{k}$, it implies that
\begin{align*}
\beta_{k}^{\frac{1}{p}}\DD_{k}f
&=\sum_{Q\in\D_{k-1}}\inner{f,h_{Q}}m(Q)^{\frac{1}{2}}\frac{\beta_{k}^{\frac{1}{p}}\chi_{Q_{-}}}{\mu(Q_{-})}-
\sum_{Q\in\D_{k-1}}\inner{f,h_{Q}}m(Q)^{\frac{1}{2}}\frac{\beta_{k}^{\frac{1}{p}}\chi_{Q_{+}}}{\mu(Q_{+})}\\
&=\sum_{Q\in\D_{k-1}}\mu(Q_{-})^{-\frac{1}{p}-1}\inner{f,h_{Q}}m(Q)^{\frac{1}{2}}\chi_{Q_{-}}-
\sum_{Q\in\D_{k-1}}\mu(Q_{+})^{-\frac{1}{p}-1}\inner{f,h_{Q}}m(Q)^{\frac{1}{2}}\chi_{Q_{+}}.
\end{align*}
For sub-cube $Q_{-},Q_{+}$ of $Q$, $Q_{-}$ and $Q_{+}$ are  disjoint.  When $p$ is a positive integer, combining the disjointness of cubes in $\D_{k-1}$, we have
\begin{align*}
|\beta_{k}^{\frac{1}{p}}\DD_{k}f|^{p}
&=\sum_{Q\in\D_{k-1}}\mu(Q_{-})^{-1-p}|\inner{f,h_{Q}}|^{p}m(Q)^{\frac{p}{2}}\chi_{Q_{-}}+
\sum_{Q\in\D_{k-1}}\mu(Q_{+})^{-1-p}|\inner{f,h_{Q}}|^{p}m(Q)^{\frac{p}{2}}\chi_{Q_{+}}.
\end{align*}
Therefore,
\begin{align*}
\Norm{\beta_{k}^{\frac{1}{p}}\DD_{k}f}_{L_{p}(\mu)}^{p}
=\sum_{Q\in\D_{k-1}}|\inner{f,h_{Q}}|^{p}m(Q)^{\frac{p}{2}}
\left(\frac{1}{\mu(Q_{-})^{p}}+\frac{1}{\mu(Q_{+})^{p}}\right).
\end{align*}
This gives the desired assertion.

\end{proof}



\begin{proposition}\label{equivalent besov}
For $0<p<\infty$,, we have
\begin{align*}
\Norm{f}_{\dot{\htb}_{p}(\mu)}\simeq
\left(\sum_{Q\in\D}|m(Q)^{-\frac{1}{2}}\inner{f,h_{Q}}|^{p}\right)^{-\frac{1}{p}},
\end{align*}
where the constants only depend on $p$.
\end{proposition}
\begin{proof}
Noting that $m(Q)=\frac{\mu(Q_{-})\mu(Q_{+})}{\mu(Q)}
\simeq\left(\frac{1}{\mu(Q_{-})},\frac{1}{\mu(Q_{+})}\right)=\min(\mu(Q_{+}),\mu(Q_{-}))$, it gives us that $$\frac{1}{\mu(Q_{-})^{p}}+\frac{1}{\mu(Q_{+})^{p}}\leq\frac{2}{m(Q)^{p}}.$$
Therefore, by \Cref{besov equi-version},
\begin{align*}
\Norm{f}_{\dot{\htb}_{p}(\mu)}\leq
2^{\frac{1}{p}}\left(\sum_{Q\in\D}|m(Q)^{-\frac{1}{2}}\inner{f,h_{Q}}|^{p}\right)^{-\frac{1}{p}}.
\end{align*}

Since $\frac{1}{m(Q)}=\frac{1}{\mu(Q_{-})}+\frac{1}{\mu(Q_{+})}$ and $(a+b)^{p}\leq\max(2^{p-1},1)(a^{p}+b^{p})$ for positive numbers $a,b$, it follows that
\begin{align*}
\frac{1}{\mu(Q_{-})^{p}}+\frac{1}{\mu(Q_{+})^{p}}
\geq\min(2^{1-p},1)(\frac{1}{\mu(Q_{-})}+\frac{1}{\mu(Q_{+})})^{p}
=\min(2^{1-p},1)\frac{1}{m(Q)^{p}}.
\end{align*}
By \Cref{besov equi-version}, we have
\begin{align*}
\Norm{f}_{\dot{\htb}_{p}(\mu)}\geq
\min(2^{\frac{1}{p}-1},1)\left(\sum_{Q\in\D}|m(Q)^{-\frac{1}{2}}\inner{f,h_{Q}}|^{p}\right)^{-\frac{1}{p}}.
\end{align*}
These imply the assertion.

\end{proof}



\begin{lemma}\label{oscillation besov}
For $0<p<\infty$, we have
\begin{align*}
\Norm{\set{f_{Q}-f_{\wh{Q}}}_{Q\in\D}}_{\ell_{p}}=\Norm{f}_{\dot{\htb}_{p}(\mu)}.
\end{align*}
\end{lemma}
\begin{proof}
For all $k\in\ZZ$,
\begin{align*}
\beta_{k}^{\frac{1}{p}}\DD_{k}f
=\sum_{Q\in\D_{k}}\beta_{k}^{\frac{1}{p}}\DD_{k}f\cdot\chi_{Q}
=\sum_{Q\in\D_{k}}\beta_{k}^{\frac{1}{p}}\left(\EE_{k}f-\EE_{k-1}f\right)\chi_{Q}
=\sum_{Q\in\D_{k}}\mu(Q)^{-\frac{1}{p}}(f_{Q}-f_{\wh{Q}})\chi_{Q},
\end{align*}
where in the last equality we use the definitions of $\beta_{k}$ and conditional expectation in \eqref{expectation definition}. 
%\begin{align*}
%\EE_{n}(f)=\sum_{Q\in\D_{n}}\frac{1}{\mu(Q)}\int_{Q}f(y)d\mu(y)\chi_{Q},\quad n\in\ZZ.
%\end{align*}
By disjointness of cubes in $\D_{k}$, we have
\begin{align*}
|\beta_{k}^{\frac{1}{p}}\DD_{k}f|^{p}
=\sum_{Q\in\D_{k}}\mu(Q)^{-1}|f_{Q}-f_{\wh{Q}}|^{p}\chi_{Q}.
\end{align*}
Moreover,
\begin{align*}
\Norm{\beta_{k}^{\frac{1}{p}}\DD_{k}f}_{L_{p}(\mu)}^{p}
=\int_{\RR}\sum_{Q\in\D_{k}}\mu(Q)^{-1}|f_{Q}-f_{\wh{Q}}|^{p}\chi_{Q}d\mu
=\sum_{Q\in\D_{k}}|f_{Q}-f_{\wh{Q}}|^{p}.
\end{align*}
Therefore,
\begin{align*}
\Norm{\set{f_{Q}-f_{\wh{Q}}}_{Q\in\D}}_{\ell_{p}}^{p}
=\sum_{Q\in\ZZ}\sum_{Q\in\D_{k}}|f_{Q}-f_{\wh{Q}}|^{p}
=\sum_{Q\in\ZZ}\Norm{\beta_{k}^{\frac{1}{p}}\DD_{k}f}_{L_{p}(\mu)}^{p}
=\Norm{f}_{\dot{\htb}_{p}(\mu)}^{p}.
\end{align*}
This is the assertion.

\end{proof}




\subsection{Schatten estimate}
For $e,h\in L_{2}(\mu)$, denote $$(e\otimes h)(f)=\inner{f,e}h,\quad f\in L_{2}(\mu).$$
Note that $(e_{1}\otimes h_{1})(e_{2}\otimes h_{2})\neq e_{1}e_{2}\otimes h_{1}h_{2}$ and $(e\otimes h)^{*}= h^{*}\otimes e^{*}$. Meanwhile, we have
\begin{align*}
{\color{red}\Norm{e\otimes h}_{S_{p}(L_{2}(\mu))}\leq\Norm{e}_{L_{2}(\mu)}\Norm{h}_{L_{2}(\mu)},
\quad 0<p<\infty.??????}
\end{align*}
Moreover, we have the following fact.
\begin{lemma}\label{tensor Sp}
Let $e_{1},e_{2},h_{1},h_{2}\in L_{2}(\mu)$ be real functions. We have
$$\Tr\left((e_{1}\otimes h_{1})^{*}(e_{2}\otimes h_{2})\right)=\inner{e_{1},e_{2}}\inner{h_{1},h_{2}}.$$
\end{lemma} 
\begin{proof}
Indeed, for $f\in L_{2}(\mu)$,
\begin{align*}
(e_{1}\otimes h_{1})^{*}(e_{2}\otimes h_{2})f
=\inner{f,e_{2}}\cdot(h_{1}\otimes e_{1}) h_{2}
=\inner{h_{2},h_{1}}\inner{f,e_{2}}e_{1}
=\inner{h_{2},h_{1}}(e_{2}\otimes e_{1})f.
\end{align*}
Since $e_{1},e_{2}\in L_{2}(\mu)$, it follows that $e_{2}\otimes e_{1}\in S_{1}(L_{2}(\mu))$. Moreover, observing that the integral kernel of $e_{2}\otimes e_{1}$ is $e_{2}(y)e_{1}(x)$, we have
\begin{align*}
\Tr(e_{2}\otimes e_{1})=\int_{\RR}e_{2}(x)e_{1}(x)d\mu(x)=\inner{e_{1},e_{2}}.
\end{align*}
Therefore,
\begin{align*}
\Tr\left((e_{1}\otimes h_{1})^{*}(e_{2}\otimes h_{2})\right)
=\inner{h_{2},h_{1}}\Tr(e_{2}\otimes e_{1})
=\inner{h_{1},h_{2}}\inner{e_{1},e_{2}}.
\end{align*}
This is the assertion.
\end{proof}


\begin{lemma}\label{paraproduct Sp}
Let $\mu$ be a positive Borel measure on $\RR$ with $0<\mu(Q)<\infty$ for $Q\in\D$. If $0<p<\infty$ and $\set{\mu(Q)^{-\frac{1}{2}}\inner{\varphi,h_{Q}}}_{Q\in\D}\in\ell_{p}$, then
$$\Norm{\pi_{\varphi}}_{S_{p}(L_{2}(\mu))}
\leq\Norm{\set{\mu(Q)^{-\frac{1}{2}}\inner{\varphi,h_{Q}}}_{Q\in\D}}_{\ell_{p}},$$
where $\pi_{\varphi}$ is given as in \Cref{L2 estimate paraproduct}.
\end{lemma}
\begin{proof}
Write
$\displaystyle\pi_{\varphi}
=\sum_{Q\in\D}\inner{\varphi,h_{Q}}\frac{\chi_{Q}}{\mu(Q)}\otimes h_{Q}.$ 
When $0<p\leq1$, it follows from triangle inequality in $\Norm{\cdot}_{S_{p}(L_{2}(\mu))}^{p}$ and orthonormality of system $\set{h_{Q}}_{Q\in\D}$ that
\begin{align*}
\Norm{\pi_{\varphi}}_{S_{p}(L_{2}(\mu))}^{p}
\leq\sum_{Q\in\D}|\inner{\varphi,h_{Q}}|^{p}
\Norm{\frac{\chi_{Q}}{\mu(Q)}\otimes h_{Q}}_{S_{p}(L_{2}(\mu))}^{p}
&\leq\sum_{Q\in\D}|\inner{\varphi,h_{Q}}|^{p}
\Norm{\frac{\chi_{Q}}{\mu(Q)}}_{L_{2}(\mu)}^{p}\Norm{h_{Q}}_{L_{2}(\mu)}^{p}\\
&=\sum_{Q\in\D}\left|\mu(Q)^{-\frac{1}{2}}\inner{\varphi,h_{Q}}\right|^{p}.
\end{align*}
This gives the assertion for $0<p\leq1$.

For $1\leq p<\infty$, let us firstly focus on $p=2k$ with $k\in\NN$. In such a case, by \Cref{tensor Sp} and using orthonormality of system $\set{h_{Q}}_{Q\in\D}$, we have
\begin{align*}
\Norm{\pi_{\varphi}}_{S_{2k}(L_{2}(\mu))}^{2k}
&=\Tr\left(\left|\sum_{Q\in\D}\inner{\varphi,h_{Q}}\frac{\chi_{Q}}{\mu(Q)}\otimes h_{Q}\right|^{2k}\right)\\
&=\Tr\left(\left(\sum_{J,Q\in\D}\inner{\varphi,h_{J}}^{*}\inner{\varphi,h_{Q}}
(\frac{\chi_{J}}{\mu(J)}\otimes h_{J})^{*}
(\frac{\chi_{Q}}{\mu(Q)}\otimes h_{Q})\right)^{k}\right)\\
&=\sum_{Q_{1},J_{1},\cdots,Q_{k},J_{k}\in\D}\inner{\varphi,h_{Q_{1}}}\inner{\varphi,h_{J_{1}}}^{*}
\cdots\inner{\varphi,h_{Q_{k}}}\inner{\varphi,h_{J_{k}}}^{*}\\
&\cdot\prod_{i=1}^{k}\inner{\frac{\chi_{Q_{i}}}{\mu(Q_{i})},\frac{\chi_{J_{i}}}{\mu(J_{i})}}
\inner{h_{J_{i}}, h_{Q_{i+1}}}\\
&=\sum_{Q_{1},\cdots,Q_{k}\in\D}\prod_{i=1}^{k}\left|\inner{\varphi,h_{Q_{i}}}\right|^{2}
\inner{\frac{\chi_{Q_{i}}}{\mu(Q_{i})},\frac{\chi_{Q_{i+1}}}{\mu(Q_{i+1})}},
\end{align*}
where we denote $Q_{k+1}=Q_{1}$. Note that $Q_{i}\subsetneq Q_{i+1}$ $(1\leq i\leq k-1)$ is impossible because the length of $Q_{i+1}$ is smaller than the length of $Q_{i}$. Therefore, due to disjointness of cubes, we can divide the index in the sum on the last line into two cases: $Q_{i}=Q_{i+1}$ and $Q_{i+1}\subsetneq Q_{i}$ $(1\leq i\leq k-1)$. When $Q_{i}=Q_{i+1}$, one has
$$\inner{\frac{\chi_{Q_{i}}}{\mu(Q_{i})},\frac{\chi_{Q_{i+1}}}{\mu(Q_{i+1})}}
=\frac{1}{\mu(Q_{i})}.$$
When $Q_{i+1}\subsetneq Q_{i}$ $(1\leq i\leq k-1)$, one has
$$\inner{\frac{\chi_{Q_{i}}}{\mu(Q_{i})},\frac{\chi_{Q_{i+1}}}{\mu(Q_{i+1})}}
=\frac{\mu(Q_{i+1})}{\mu(Q_{i})\mu(Q_{i+1})}
=\frac{1}{\mu(Q_{i})}.$$
This means that we have the same result whenever in the above which cases. Thus, by property of combinatorics, we arrive at
\begin{align*}
\Norm{\pi_{\varphi}}_{S_{2k}(L_{2}(\mu))}^{2k}
=\sum_{Q_{1},\cdots,Q_{k}\in\D}\prod_{i=1}^{k}\left|\inner{\varphi,h_{Q_{i}}}\right|^{2}
\inner{\frac{\chi_{Q_{i}}}{\mu(Q_{i})},\frac{\chi_{Q_{i+1}}}{\mu(Q_{i+1})}}
=k2^{k-1}\sum_{Q\in\D}\left|\mu(Q)^{-\frac{1}{2}}\inner{\varphi,h_{Q}}\right|^{2k}.
\end{align*}
Moreover, by interpolation ???,


...............................
\end{proof}




\begin{lemma}\label{ONB Sp}
Let $\mu$ is a positive Borel measure on $\RR$ with $0<\mu(Q)<\infty$ for $Q\in\D$. For orthonormal functions $\set{e_{Q}}_{Q\in\D}$  and scalar sequence $\set{\lambda_{Q}}_{Q\in\D}$, define $$\C_{\lambda}=\sum_{Q\in\D}\lambda_{Q}e_{Q}\otimes h_{Q}.$$
If $0<p<\infty$ and $\set{\lambda_{Q}}_{Q\in\D}\in\ell_{p}$, then 
$$\Norm{\C_{\lambda}}_{S_{p}(L_{2}(\mu))}\lesssim\Norm{\set{\lambda_{Q}}_{Q\in\D}}_{\ell_{p}}.$$
\end{lemma}
\begin{proof}
This proof is similar to \Cref{paraproduct Sp} but more easy.
When $0<p\leq1$, it follows from triangle inequality in $\Norm{\cdot}_{S_{p}(L_{2}(\mu))}^{p}$, and orthonormality of system $\set{e_{Q}}_{Q\in\D}$ and $\set{h_{Q}}_{Q\in\D}$ that
\begin{align*}
\Norm{\C_{\lambda}}_{S_{p}(L_{2}(\mu))}^{p}
\leq\sum_{Q\in\D}|\lambda_{Q}|^{p}
\Norm{e_{Q}\otimes h_{Q}}_{S_{p}(L_{2}(\mu))}^{p}
\leq\sum_{Q\in\D}|\lambda_{Q}|^{p}
\Norm{e_{Q}}_{L_{2}(\mu)}^{p}\Norm{h_{Q}}_{L_{2}(\mu)}^{p}
=\sum_{Q\in\D}\left|\lambda_{Q}\right|^{p}.
\end{align*}
This yields the assertion of $0<p\leq1$.
For $1\leq p<\infty$, let us firstly focus on $p=2k$ with $k\in\NN$. In such a case, by \Cref{tensor Sp} and using orthonormality of system $\set{e_{Q}}_{Q\in\D}$ and $\set{h_{Q}}_{Q\in\D}$, we have
\begin{align*}
\Norm{\C_{\lambda}}_{S_{2k}(L_{2}(\mu))}^{2k}
&=\Tr\left(\left|\sum_{Q\in\D}\lambda_{Q}e_{Q}\otimes h_{Q}\right|^{2k}\right)
=\Tr\left(\left(\sum_{J,Q\in\D}\lambda_{J}^{*}\lambda_{Q}
(e_{J}\otimes h_{J})^{*}(e_{Q}\otimes h_{Q})\right)^{k}\right)\\
&=\sum_{Q_{1},J_{1},\cdots,Q_{k},J_{k}\in\D}\lambda_{Q_{1}}\lambda_{J_{1}}^{*}
\cdots\lambda_{Q_{k}}\lambda_{J_{k}}^{*}\cdot\prod_{i=1}^{k}\inner{e_{Q_{i}},e_{J_{i}}}
\inner{h_{J_{i}}, h_{Q_{i+1}}}\\
&=\sum_{Q_{1},J_{1},\cdots,Q_{k},J_{k}\in\D}\lambda_{Q_{1}}\lambda_{J_{1}}^{*}
\cdots\lambda_{Q_{k}}\lambda_{J_{k}}^{*}\cdot\prod_{i=1}^{k}\delta_{Q_{i},J_{i}}\delta_{J_{i},Q_{i+1}}\\
&=\sum_{Q\in\D}\left|\lambda_{Q}\right|^{2k},
\end{align*}
where we denote $Q_{k+1}=Q_{1}$ and $\delta_{Q,J}$ is the $\delta$-function. Therefore, by interpolation ???,





...............................
\end{proof}



\begin{corollary}\label{L0 Sp}
Let $\mu$ is a positive Borel measure on $\RR$ with $0<\mu(Q)<\infty$ for $Q\in\D$. If $0<p<\infty$ and $\varphi\in\dot{\htb}_{p}(\mu)$, then 
$$\Norm{\Lambda^{0}_{\varphi}}_{S_{p}(L_{2}(\mu))}
\lesssim\Norm{\varphi}_{\dot{\htb}_{p}(\mu)}.$$
where $\Lambda^{0}_{\varphi}$ is given as in \Cref{L2 estimate ad-paraproduct}.
\end{corollary}
\begin{proof}
Since $\DD_{Q}$ here is a self-adjoint operator, it yields that
$\displaystyle\Lambda^{0}_{\varphi}=\sum_{Q\in\D}\DD_{Q}\left(h_{Q}\cdot\DD_{Q}\varphi\right)\otimes h_{Q}$.
Combining definition of $\DD_{Q}$, we have
$$\DD_{Q}\left(h_{Q}\cdot\DD_{Q}\varphi\right)=\inner{\varphi,h_{Q}}\inner{h_{Q}^{3},1}h_{Q}.$$
Therefore,
\begin{align*}
\Lambda^{0}_{\varphi}=
\sum_{Q\in\D}\inner{\varphi,h_{Q}}\inner{h_{Q}^{3},1}h_{Q}\otimes h_{Q}.
\end{align*}
By \Cref{ONB Sp}, $\Norm{h_{Q}}_{L_{3}(\mu)}\simeq m(Q)^{-\frac{1}{6}}$ and \Cref{equivalent besov}, we conclude that
\begin{align*}
\Norm{\Lambda^{0}_{\varphi}}_{S_{p}(L_{2}(\mu))}
\lesssim\left(\sum_{Q\in\D}|\inner{\varphi,h_{Q}}\inner{h_{Q}^{3},1}|^{p}\right)^{\frac{1}{p}}
\lesssim\left(\sum_{Q\in\D}|m(Q)^{-\frac{1}{2}}\inner{\varphi,h_{Q}}|^{p}\right)^{\frac{1}{p}}
\lesssim\Norm{\varphi}_{\dot{\htb}_{p}(\mu)}.
%&\lesssim\left(\sum_{Q\in\D,\mu(Q_{+})\leq\mu(Q_{-})}|m(Q)^{-\frac{1}{2}}\inner{\varphi,h_{Q}}|^{p}\right)^{\frac{1}{p}}
%+\left(\sum_{Q\in\D,\mu(Q_{-})\leq\mu(Q_{+})}|m(Q)^{-\frac{1}{2}}\inner{\varphi,h_{Q}}|^{p}\right)^{\frac{1}{p}}\\
%&\lesssim\Norm{\set{\mu(Q)^{-\frac{1}{2}}\inner{\varphi,h_{\wh{Q}}}}_{Q\in\D}}_{\ell_{p}}.
\end{align*}
This is the assertion.


\end{proof}





\begin{proposition}\label{commutator Sp estimate}
Let $\mu$ be a positive Borel measure on $\RR$ with $0<\mu(Q)<\infty$ for each $Q\in\D$.  If $\varphi\in\dot{\htb}_{p}(\mu)$, then 
\begin{align*}
\Norm{[M_{\varphi},H_{\D}]}_{S_{p}(L_{2}(\mu))}
\leq C\Norm{\varphi}_{\dot{\htb}_{p}(\mu)},
\end{align*}
for some positive constant $C$ depending only on $p$.
\end{proposition}
\begin{proof}
As in the proof of \Cref{L2 estimate}, take the following decomposition into a consideration
\begin{align*}
M_{\varphi}=\pi_{\varphi}+\Lambda^{0}_{\varphi}+\Lambda^{1}_{\varphi}+\pi_{\varphi}^{*},
\end{align*}
where
\begin{align*}
\pi_{\varphi}f
=\sum_{Q\in\D}\DD_{Q}\varphi\cdot \EE_{Q}f,\quad
\pi_{\varphi}^{*}\; \mbox{is\; the}\; L_{2}(\mu)\mbox{-adjoint\; of}\; \pi_{\varphi},
\end{align*}
and
\begin{align*}
\Lambda^{0}_{\varphi}f
=\sum_{Q\in\D}\DD_{Q}\left(\DD_{Q}\varphi\cdot\DD_{Q}f\right),\quad
\Lambda^{1}_{\varphi}f
=\sum_{Q\in\D}\EE_{Q}\varphi\cdot\DD_{Q}f.
\end{align*}
Therefore,
$$[M_{\varphi},H_{\D}]=[\pi_{\varphi},H_{\D}]+[\Lambda^{0}_{\varphi},H_{\D}]
+[\Lambda^{1}_{\varphi},H_{\D}]+[\pi_{\varphi}^{*},H_{\D}].$$ 
It suffices to investigate the Schatten estimates of $[\pi_{\varphi},H_{\D}]$, $[\pi_{\varphi}^{*},H_{\D}]$, $[\Lambda^{0}_{\varphi},H_{\D}]$ and $[\Lambda^{1}_{\varphi},H_{\D}]$.

Note that the operator $H_{\D}$ is bounded on $L_{2}(\mu)$ with upper bound $\sqrt{2}$ when $\mu$ is a positive Borel measure. Since, for $Q\in\D$,
\begin{align*}
\frac{1}{\mu(Q)}\leq\min(\frac{1}{\mu(Q_{+})},\frac{1}{\mu(Q_{-})})
\leq\frac{1}{\min(\mu(Q_{+}),\mu(Q_{+}))}=\frac{1}{m(Q)},
\end{align*}
it follows from \Cref{paraproduct Sp} and \Cref{equivalent besov} that
\begin{align*}
\Norm{[\pi_{\varphi},H_{\D}]}_{S_{p}(L_{2}(\mu))}
%\leq\Norm{\pi_{\varphi}H_{\D}}_{S_{p}(L_{2}(\mu))}+\Norm{H_{\D}\pi_{\varphi}}_{S_{p}(L_{2}(\mu))}
\leq2\Norm{\pi_{\varphi}}_{S_{p}(L_{2}(\mu))}\Norm{H_{\D}}_{\B(L_{2}(\mu))}
&\leq2\sqrt{2}\Norm{\set{\mu(Q)^{-\frac{1}{2}}\inner{\varphi,h_{Q}}}_{Q\in\D}}_{\ell_{p}}\\
&\lesssim\Norm{\set{m(Q)^{-\frac{1}{2}}\inner{\varphi,h_{Q}}}_{Q\in\D}}_{\ell_{p}}\\
&\lesssim\Norm{\varphi}_{\dot{\htb}_{p}(\mu)}.
\end{align*} 
At the same time, since $H_{\D}^{*}$ is bounded on $L_{2}(\mu)$ with upper bound $\sqrt{2}$ and $\Norm{\cdot}_{S_{p}(L_{2}(\mu))}$ is unitarily invariant,  it follows from \Cref{paraproduct Sp} and \Cref{equivalent besov} that
\begin{align*}
\Norm{[\pi_{\varphi}^{*},H_{\D}]}_{S_{p}(L_{2}(\mu))}
=\Norm{[H_{\D}^{*},\pi_{\varphi}]^{*}}_{S_{p}(L_{2}(\mu))}
&=\Norm{[H_{\D}^{*},\pi_{\varphi}]}_{S_{p}(L_{2}(\mu))}\\
&\leq2\Norm{\pi_{\varphi}}_{S_{p}(L_{2}(\mu))}\Norm{H_{\D}^{*}}_{\B(L_{2}(\mu))}\\
&\leq2\sqrt{2}\Norm{\set{\mu(Q)^{-\frac{1}{2}}\inner{\varphi,h_{Q}}}_{Q\in\D}}_{\ell_{p}}\\
&\lesssim\Norm{\set{m(Q)^{-\frac{1}{2}}\inner{\varphi,h_{Q}}}_{Q\in\D}}_{\ell_{p}}\\
&\lesssim\Norm{\varphi}_{\dot{\htb}_{p}(\mu)}.
\end{align*} 
For $[\Lambda^{0}_{\varphi},H_{\D}]$, by \Cref{L0 Sp},
\begin{align*}
\Norm{[\Lambda^{0}_{\varphi},H_{\D}]}_{S_{p}(L_{2}(\mu))}
\leq2\Norm{\Lambda^{0}_{\varphi}}_{S_{p}(L_{2}(\mu))}\Norm{H_{\D}}_{\B(L_{2}(\mu))}
\lesssim\Norm{\varphi}_{\dot{\htb}_{p}(\mu)}.
\end{align*} 
To investigate the Schatten estimate of $[\Lambda^{1}_{\varphi},H_{\D}]$, let us revisit the proof of \Cref{L2 estimate}. It gives that 
\begin{align*}
[\Lambda^{1}_{\varphi},H_{\D}]=
\sum_{Q\in\D}\sigma_{Q}\left(\varphi_{Q}-\varphi_{\wh{Q}}\right)h_{\wh{Q}}\otimes h_{Q}.
\end{align*}
Since $\set{h_{\wh{Q}}}$ are orthonormal functions, by \Cref{ONB Sp} and \Cref{oscillation besov}, one has
\begin{align*}
\Norm{[\Lambda^{1}_{\varphi},H_{\D}]}_{S_{p}(L_{2}(\mu))}
\lesssim\Norm{\set{\varphi_{Q}-\varphi_{\wh{Q}}}_{Q\in\D}}_{\ell_{p}}
\lesssim\Norm{\varphi}_{\dot{\htb}_{p}(\mu)}.
\end{align*}
Taking the estimates on $[\pi_{\varphi},H_{\D}]$, $[\pi_{\varphi}^{*},H_{\D}]$, $[\Lambda^{0}_{\varphi},H_{\D}]$ and $[\Lambda^{1}_{\varphi},H_{\D}]$ in a combination, we claim the assertion.

\end{proof}









{\bf Acknowledgments:} F. Sukochev and D. Zanin are supported by ARC DP230100434. 
F. L. Yang is supported by National Natural Science Foundation of China No. 12371138 and he is grateful for the support of his supervisor professor Xiao Xiong.




\bibliographystyle{amsplain}
\bibliography{NCZhou}




%\begin{thebibliography}{99}

%\bibitem{CRW1976}R.R. Coifman, R. Rochberg, G. Weiss. 
%{\it Factorization theorems for Hardy spaces in several variables.} 
%Ann. of Math. 103 (1976) 611--635.

%\bibitem{Pe1995}C. P\'{e}rez. 
%{\it Endpoint estimates for commutators of singular integral operators.} 
%J. Funct. Anal., 128 (1995), 163--185.



%\bibitem{LMP2014}López-Sánchez, Luis Daniel; Martell, José María; Parcet, Javier. 
%{\it Dyadic harmonic analysis beyond doubling measures.} 
%Adv. Math. 267 (2014), 44--93.
%
%\bibitem{PKJF-book}L. Pick, A. Kufner, O. John and S. Fu\v{c}ík. 
%{\it Function spaces, 1.} 
%Walter de Gruyter, 2012.



%\bibitem{Treil2013}Treil, Sergei. 
%{\it Commutators, paraproducts and BMO in non-homogeneous martingale settings.}
%Rev. Mat. Iberoam. 29 (2013), no. 4, 132--1372.

%\bibitem{Gr-book1}L. Grafakos. 
%{\it Classical Fourier analysis.} 
%Third edition. Graduate Texts in Mathematics, 249. Springer, New York, 2014. xviii+638 pp.

%\bibitem{G2014m}Grafakos, Loukas. 
%{\it Modern Fourier analysis.} 
%Third edition. Graduate Texts in Mathematics, 250. Springer, New York, 2014. xvi+624 pp.

%\bibitem{P2000}S. Petermichl. 
%{\it Dyadic shifts and a logarithmic estimate for Hankel operators with matrix symbol.}
%C. R. Acad. Sci. Paris Sér. I Math. 330 (6) (2000) 455--460.

%\bibitem{CMP2012}D. V. Cruz-Uribe, J. M. Martell, C. Pérez. 
%{\it Sharp weighted estimates for classical operators.} 
%Adv. Math. 229 (1) (2012) 408--441.

%\bibitem{Weisz-book1994}F. Weisz. 
%{\it Martingale Hardy Spaces and Their Applications in Fourier Analysis.} 
%Lecture Notes in Mathematics, vol. 1568, Springer-Verlag, Berlin, 1994.

%\bibitem{G1968}R.F. Gundy. 
%{\it A decomposition for L1-bounded martingales.} 
%Ann. Math. Statist. 39 (1968) 134--138.

%\bibitem{CMP-book2011}D.V. Cruz-Uribe, J.M. Martell, C. Pérez. 
%{\it Weights, Extrapolation and the Theory of Rubio de Francia.} 
%Oper. Theory Adv. Appl., vol. 215, Birkhäuser/Springer Basel AG, Basel, 2011.

%\bibitem{M2003}Mei, Tao. 
%{\it BMO is the intersection of two translates of dyadic BMO.} 
%C. R. Math. Acad. Sci. Paris 336 (2003), no. 12, 1003--1006


%\end{thebibliography}




\end{document}
